\section{Preliminaries}\label{prel}

\subsection{Reduced affine root systems}

Let $(E,\langle\cdot,\cdot\rangle)$ be an Euclidean space of dimension $r$. Transferring the inner product $\langle\cdot,\cdot\rangle$ on~$E$ to~$E^*$ through the linear isomorphism \smash{$E\overset{\sim}{\longrightarrow} E^*$}, $y\mapsto \langle y,\cdot\rangle$, is turning $E^*$ into an Euclidean space. We denote its inner product again by $\langle\cdot,\cdot\rangle$, and its norm by $\|\cdot\|$.

Let $\Phi_0$ be an irreducible reduced root system in $E^*$ with Weyl group $W_0$. Its dual root system $\Phi_0^\vee=\{\alpha^\vee\}_{\alpha\in\Phi_0}$ in $E$ consists of the co-roots $\alpha^\vee\in E$ ($\alpha\in\Phi_0$), which are the vectors in~$E$ satisfying
\begin{equation}\label{corootdef}
\bigl\langle y,\alpha^\vee\bigr\rangle=\frac{2\alpha(y)}{\|\alpha\|^2}
\end{equation}
for all $y\in E$.

Consider the corresponding reduced affine root system
\[
\Phi:=\Phi_0\times \mathbb{Z}\subset E^*\times\mathbb{R}.
\]
We will view an element $(\phi,\xi)\in E^*\times\mathbb{R}$ in the ambient space as an affine linear functional on~$E$ by $y\mapsto\phi(y)+\xi$ for $y\in E$.

The projection $E^*\times\mathbb{R}\rightarrow E^*$ on the first component will be denoted by
\[
f\mapsto\overline{f}.
\]
It restricts to a surjective map $\Phi\twoheadrightarrow\Phi_0$. Furthermore, we have $a=(\overline{a},a(0))$ for $a\in\Phi$.
Throughout the paper, we will identify a root $\alpha\in\Phi_0$ with $(\alpha,0)\in\Phi$.

For $a\in\Phi$, denote by $s_a\colon E\rightarrow E$ the orthogonal reflection in the affine root hyperplane $a^{-1}(0)\subset E$. Then
\begin{equation*}%\label{reflformula}
s_a(y)=y-a(y)\overline{a}^\vee
\end{equation*}
for $y\in E$.
The affine Weyl group $W$ of $\Phi$ is the subgroup of affine linear transformations of $E$ generated by the orthogonal reflections $s_a$, $a\in\Phi$. The finite Weyl group $W_0$ is the subgroup generated by $s_\alpha$, $\alpha\in\Phi_0$.

For $y\in E$, let $\tau(y)\colon E\rightarrow E$ be the translation map $z\mapsto z+y$. Then
\begin{equation}\label{ssplit}
s_a=s_{\overline{a}}\,\tau\bigl(a(0)\overline{a}^\vee\bigr)
\end{equation}
for $a\in\Phi$. Consequently,
$W\simeq W_0\ltimes Q^\vee$ with $Q^\vee=\mathbb{Z}\Phi_0^\vee$ the co-root lattice of $\Phi_0$.

The linear, contragredient $W$-action
on the space $E^*\times\mathbb{R}$ of affine linear functionals on $E$ restricts to a $W$-action on $\Phi$. It satisfies
\begin{align}
&s_a(b)=b-\overline{b}\bigl(\overline{a}^\vee\bigr)a=\bigl(s_{\overline{a}}(\overline{b}),b(0)-a(0)\overline{b}\bigl(\overline{a}^\vee\bigr)\bigr),\nonumber\\
&\tau(\lambda)b=\bigl(\overline{b},b(0)-\overline{b}(\lambda)\bigr)\label{WPhiaction}
\end{align}
for $a,b\in\Phi$ and $\lambda\in Q^\vee$.
%%%%%%%%%%%%%%%%%%%%%%%%%%%%%%%%%%%%%%%%

We fix an ordered basis $\Delta_0=\{\alpha_1,\dots,\alpha_r\}$ of the root system $\Phi_0$ once and for all. We will choose the ordering such that the following convention
holds true.

\begin{Convention}\label{root_length_convention}
The simple root $\alpha_r$ is a long root.
\end{Convention}

If all the roots in $\Phi_0$ have the same root length, then all roots are considered to be long as well as short.
We denote by $\Phi_0^{+}$ the set of positive roots in $\Phi_0$ relative to $\Delta_0$. The corresponding set of negative roots is denoted by $\Phi_0^-:=-\Phi_0^+$.

The Weyl group $W_0$ is a Coxeter group with Coxeter generators $\{s_1,\dots,s_r\}$ given by the simple reflections $s_i:=s_{\alpha_i}$, $1\leq i\leq r$. The closure of the positive Weyl chamber
\begin{equation*}%\label{Eplus}
E_+:=\bigl\{y\in E \mid \alpha(y)>0\ \forall \alpha\in\Phi_0^+ \bigr\}
\end{equation*}
is a fundamental domain for the $W_0$-action on $E$.

The ordered basis $\Delta_0$ of $\Phi_0$ extends to an ordered basis
$\Delta=\{\alpha_0,\dots,\alpha_r\}$ of $\Phi$ with the additional affine simple root
\[
\alpha_0=(-\varphi,1),
\]
where $\varphi$ is the highest root of $\Phi_0$ relative to $\Delta_0$. The corresponding sets of positive and negative roots are denoted by $\Phi^+$ and $\Phi^-$, respectively.
The affine Weyl group $W$ is a Coxeter group with Coxeter generators $\{s_0,\dots,s_r\}$ the simple reflections $s_j:=s_{\alpha_j}$, $0\leq j\leq r$. By~\eqref{ssplit}, we~have
\[
s_0=s_\varphi\tau\bigl(-\varphi^\vee\bigr)=\tau\bigl(\varphi^\vee\bigr)s_\varphi.
\]

The closure $\overline{C_+}$ of the fundamental alcove
\begin{equation*}%\label{Cplus}
C_+:=\{y\in E_+ \mid \alpha_0(y)>0\}
\end{equation*}
is a fundamental domain for the action of $W\simeq W_0\ltimes Q^\vee$ on $E$ by reflections and translations.
For a $W$-orbit $\mathcal{O}$ in $E$, we denote by
$c^\mathcal{O}$ the unique vector in $\mathcal{O}\cap \overline{C_+}$.

Since $\alpha_0=(-\varphi,1)$, we have the following alternative description:
\[
C_+=\bigl\{y\in E \mid 0<\alpha(y)<1 \ \forall \alpha\in\Phi_0^+\bigr\}
\]
of the fundamental alcove.
Note furthermore that
\begin{equation}\label{W0C}
\bigcup_{w\in W_0}w\bigl(\overline{C_+}\bigr)=\{y\in E \mid |\alpha(y)|\leq 1\ \forall \alpha\in\Phi_0\}.
\end{equation}

\subsection{Nonreduced extensions and multiplicity functions}

For $a\in\Phi$ such that $\overline{a}\smash{\bigl(Q^\vee\bigr)}=\mathbb{Z}$, we have
\[
Wa=W_0 \overline{a}\times\mathbb{Z}
\]
by~\eqref{WPhiaction}. A case by case inspection of the Dynkin diagrams shows that $\alpha\smash{\bigl(Q^\vee\bigr)}=\mathbb{Z}$ for $\alpha\in\Phi_0$ unless $\alpha\in\Phi_0$ is long and
 $\Phi_0$ is of type ${\rm C}_r$, $r\geq 1$, in which case $\alpha\smash{\bigl(Q^\vee\bigr)}=2\mathbb{Z}$ (note that ${\rm C}_1=\textup{A}_1$ and ${\rm C}_2=\textup{B}_2$).
If $\Phi_0$ is of type ${\rm C}_r$, $r\geq 1$, then $\alpha_r$ is the only long simple root in $\Delta_0$ in view of convention~\ref{root_length_convention}.

The $W$-orbits in $\Phi$ can now be described as follows:
\begin{enumerate}\itemsep=0pt
\item[$(1)$] If all the roots in $\Phi_0$ have the same root length but $\Phi_0$ is not of type $\textup{A}_1$, then $W$ acts transitively on $\Phi$.
\item[$(2)$] If $\Phi_0$ is of type ${\rm C}_1=\textup{A}_1$, then $\Phi=W\alpha_0\sqcup W\alpha_1$ and
\[
W\alpha_0=\Phi_0\times\mathbb{Z}_o,\qquad W\alpha_1=\Phi_0\times\mathbb{Z}_e
\]
with $\mathbb{Z}_o$ (resp.\ $\mathbb{Z}_e$) the set of odd (resp.\ even) integers.
\item[$(3)$] If $\Phi_0$ is of type ${\rm C}_r$, $r\geq 2$, then $\Phi=W\alpha_0\sqcup W\alpha_1\sqcup W\alpha_r$ and $\alpha_i\in W\alpha_1$ for all $1\leq i<r$. Furthermore,
\[
W\alpha_0=\Phi_0^\ell\times\mathbb{Z}_o,\qquad W\alpha_1=\Phi_0^s\times\mathbb{Z},\qquad W\alpha_r=\Phi_0^\ell\times\mathbb{Z}_e
\]
with $\Phi_0^\ell$ (resp.\ $\Phi_0^s$) the long (resp.\ short) roots in $\Phi_0$.
\item[$(4)$] If $\Phi_0$ is of type $\textup{B}_r$, $r\geq 3$, $F_4$ or $G_2$ and if $\alpha_i\in\Delta_0$, $1\leq i<r$, is a short simple root, then $\Phi=W\alpha_0\sqcup W\alpha_i$ and
\[
W\alpha_0=\Phi_0^\ell\times\mathbb{Z}=W\alpha_r,\qquad W\alpha_i=\Phi_0^s\times\mathbb{Z}.
\]
\end{enumerate}
The set
\[
\Phi^{\textup{nr}}:=\Phi\sqcup\bigl\{a/2\mid a\in\Phi\ \textup{such that}\ \overline{a}\bigl(Q^\vee\bigr)=\mathbb{Z}_e\big\}
\]
forms an affine root system in the affine space $E^*\times\mathbb{R}$ (see~\cite{MacAff}). If $\Phi_0$ is of type $C_r$, $r\geq 1$, then $\Phi^{\textup{nr}}$ is the nonreduced irreducible affine root system of type ${\rm C}^\vee{\rm C}_r$ (see~\cite{MacAff}). In this case, $\Phi^{\textup{nr}}$~has five $W$-orbits $W\alpha_0$, $W\frac{\alpha_0}{2}$, $W\alpha_1$, $W\alpha_r$, $W\frac{\alpha_r}{2}$ when $r\geq 2$, and four $W$-orbits
$W\alpha_0$, $W\frac{\alpha_0}{2}$, $W\alpha_1$, $W\frac{\alpha_1}{2}$ when $r=1$. If $\Phi_0$ is not of type ${\rm C}_r$, $r\geq 1$, then $\Phi^{\textup{nr}}=\Phi$.

Let $\mathbf{F}$ be a field of characteristic zero. We call a $W$-invariant function
\[
\mathbf{k}\colon\ \Phi^{\textup{nr}}\rightarrow\mathbf{F}^\times,\qquad a\mapsto \mathbf{k}_a
\]
a multiplicity function. Denote by $\mathcal{K}$ the set of multiplicity functions.
In order to obtain uniform notations, we extend a multiplicity function
$\mathbf{k}\colon \Phi^{\textup{nr}}\rightarrow\mathbf{F}^\times$ to a $W$-invariant function
\[
\mathbf{k}\colon\ \Phi\sqcup\tfrac{1}{2}\Phi\rightarrow\mathbf{F}^\times
\]
by declaring \smash{$\mathbf{k}_{\frac{a}{2}}:=\mathbf{k}_{a}$} when \smash{$\frac{a}{2}\not\in\Phi^{\textup{nr}}$}. Note that for any root $\alpha\in\Phi_0$,
\begin{equation}
\label{equalparameter}
\mathbf{k}_\alpha=\mathbf{k}_{(\alpha,1)}=\mathbf{k}_{\frac{\alpha}{2}}=\mathbf{k}_{(\frac{\alpha}{2},\frac{1}{2})}\qquad \textup{if $\Phi_0$ is not of type ${\rm C}_r$, $r\geq 1$}.
\end{equation}
For the value of a multiplicity function $\mathbf{k}$ at a simple root $\alpha_j$ and at \smash{$\frac{\alpha_j}{2}$}, we will use the shorthand notations
\[
k_j:=\mathbf{k}_{\alpha_j},\qquad u_j:=\mathbf{k}_{\frac{\alpha_j}{2}}.
\]

If $\Phi_0$ has rank $r=1$, then
the multiplicity function $\mathbf{k}$ is determined by the four parameters $k_0$, $u_0$, $k_1$, $u_1$, which can be chosen arbitrarily.
If $\Phi_0$ has rank $r>1$, then $\mathbf{k}$ is determined by the five parameters $k_0$, $u_0$, $k:=k_i$, $k_r$, $u_r$, where $1\leq i<r$ is such that $\alpha_i$ is a short root. These parameters can be chosen arbitrarily
when $\Phi_0$ is of type~${\rm C}_r$, $r\geq 2$. If $\Phi_0$ is not of type~${\rm C}_r$, $r\geq 1$,
then $k_0=u_0=k_r=u_r$ by~\eqref{equalparameter}, hence $\mathcal{K}\simeq\mathbf{F}^\times$ if all the roots in $\Phi_0$ have the same length and $\mathcal{K}\simeq (\mathbf{F}^\times)^2$ otherwise (i.e., if $\Phi_0$ is of type $\textup{B}_r$, $r\geq 3$, $F_4$ or $G_2$).

The extended affine Weyl group is the subgroup
\[
W^{\textup{ext}}:=W_0\tau\bigl(P^\vee\bigr)
\]
of the group of affine linear transformations of $E$,
with $P^\vee$ the co-weight lattice of $\Phi_0$. The linear, contragredient $W^{\textup{ext}}$-action on the space $E^*\times\mathbb{R}$ of affine linear functionals on $E$ restricts to a $W^{\textup{ext}}$-action on $\Phi$. The explicit formulas for this action are again given by~\eqref{WPhiaction}, now with $\lambda\in P^\vee$ in the second formula.

Write $\mathcal{K}^{{\rm res}}\subseteq\mathcal{K}$ for the subset of multiplicity functions $\mathbf{k}$ satisfying the following
two additional conditions:
\begin{enumerate}\itemsep=0pt
\item[$ (1)$] $\mathbf{k}$ is $W^{\textup{ext}}$-invariant,
\item[$(2)$] $\mathbf{k}_{\frac{a}{2}}=\mathbf{k}_a$ for all $a\in\Phi$.
\end{enumerate}
By $(2)$, a restricted multiplicity function $\mathbf{k}\in\mathcal{K}^{{\rm res}}$ is uniquely determined by its values on $\Phi$, and $\mathbf{k}_a=\mathbf{k}_{\overline{a}}$ for $a\in\Phi$. Its value $\mathbf{k}_\alpha$ at $\alpha\in\Phi_0$ only depends on the length of $\alpha$. Hence $\mathcal{K}^{{\rm res}}\simeq\mathbf{F}^\times$ if all roots in $\Phi_0$ have the same length and
$\mathcal{K}^{{\rm res}}\simeq (\mathbf{F}^\times)^2$ otherwise, which implies that
\[
\mathcal{K}^{{\rm res}}=\mathcal{K}\qquad \textup{if }  \Phi_0 \ \textup{is not of type}\ {\rm C}_r,\ r\geq 1.
\]
In particular, for $\mathbf{k}\in\mathcal{K}^{{\rm res}}$ formula~\eqref{equalparameter} holds true for root systems $\Phi_0$ of any type,
\[
k_0=u_0=k_r=u_r\qquad \text{if}\ \mathbf{k}\in\mathcal{K}^{{\rm res}}.
\]

Define an involution
\begin{equation*}%\label{dualityinvolution}
\mathcal{K}\overset{\sim}{\longrightarrow}\mathcal{K},\qquad \mathbf{k}\mapsto \widetilde{\mathbf{k}}
\end{equation*}
by interchanging the values $k_0$ and $u_r$ of $\mathbf{k}$ on the $W$-orbits $W\alpha_0$ and \smash{$W\frac{\alpha_r}{2}$}. We call it the duality involution.
Note that it is the identity unless $\Phi_0$ is of type ${\rm C}_r$, $r\geq 1$. Furthermore, its restriction to $\mathcal{K}^{{\rm res}}$ is the identity for all types.

The standard Cherednik--Macdonald theory for parameters $\mathbf{k}\in\mathcal{K}^{{\rm res}}$ admits an extension to
parameters $\mathbf{k}\in\mathcal{K}$~\cite{No,Sa}. It only gives new results when $\Phi_0$ is of type ${\rm C}_r$, $r\geq 1$, since otherwise $\mathcal{K}=\mathcal{K}^{{\rm res}}$.
The resulting theory is sometimes referred to as the Koornwinder case \big(since the associated analogs of the symmetric Macdonald polynomials are the Koornwinder polynomials~\cite{Ko}\big), or as the ${\rm C}^\vee{\rm C}_r$ case (since $\mathcal{K}$ is the natural set of multiplicity functions on the nonreduced root system of type
${\rm C}^\vee{\rm C}_r$).
It is common in the literature on Koornwinder~$\big({\rm C}^\vee{\rm C}_r\big)$ extensions of the Cherednik--Macdonald theory
to develop the theory directly using the following explicit realisation of the root system $\Phi_0$ of type ${\rm C}_r$, $r\geq 1$, see, e.g.,~\cite{Ko,No,Sa,St}:
\begin{itemize}\itemsep=0pt
\item $E=\mathbb{R}^r$ with orthonormal basis $\{\epsilon_i\}_{i=1}^r$,
\item $\Phi_0^s=\{\pm(\epsilon_i\pm\epsilon_j)\}_{1\leq i<j\leq r}$ ($=\varnothing$ when $r=1$) and $\Phi_0^\ell=\{\pm 2\epsilon_i\}_{i=1}^r$, where we identify~${E\simeq E^*}$
via the scalar product \big(in particular, $Q^\vee=\bigoplus_{i=1}^r\mathbb{Z}\epsilon_i$\big),
\item $\alpha_i=\epsilon_i-\epsilon_{i+1}$, $1\leq i<r$, and $\alpha_r=2\epsilon_r$.
\end{itemize}
Note that
$\varphi=2\epsilon_1$ is the highest root, hence $\alpha_0=(-2\epsilon_1,1)$. With these choices, the type~${\rm C}_r$ results presented in Sections~\ref{1.2}--\ref{1.4} follow immediately from the general results as discussed below. In the remainder of the paper, $\mathbf{k}$ will be a multiplicity function in $\mathcal{K}$ unless stated explicitly otherwise.

\begin{Remark}
The setup in this subsection follows~\cite{StAM, StCh}. In terms of the initial data $D$ from~\cite[Section~1.1]{StAM}, we are considering the case of twisted adjoint root data $D=(R_0,t,\Lambda,\Lambda)$ with $\Lambda$ the root lattice of $R_0$. Then
$(\Phi_0,\Phi^{\textup{nr}})$ corresponds to $\bigl(R_0^\vee,R(D)^\vee\bigr)$, with $R(D)^\vee$ the dual of the (possibly non-reduced) affine root system $R(D)$ from~\cite[Section~1.1]{StAM}. All nonreduced cases of the Cherednik--Macdonald theory as described in Macdonald's book~\cite{Ma} can be recovered from the case that $\Phi_0$ is of type ${\rm C}_r$ by appropriate specialisations of the multiplicity parameters, see, e.g.,~\cite[Section~9.2.3]{StCh} for details.
\end{Remark}
%%%%%%%%%%%%%%%%%%%%%%%%%%%%%%%%%%%%%%%%%%%%%%



%%%%%%%%%%%%%%%%%%%%%%%%%%%%%%%%%%%%
\subsection{Quasi-polynomials}%\label{qpsection}
%%%%%%%%%%%%%%%%%%%%%%%%%%%%%%%%%%%%
The space of quasi-polynomials~\cite{SSV2} is the group algebra
\[
\mathbf{F}[E]=\bigoplus_{y\in E}\mathbf{F}x^y
\]
of $E$, viewed as abelian additive group. Here we denote the canonical basis elements $x^y$, ${y\in E}$, of~$\mathbf{F}[E]$ multiplicatively, so \smash{$x^y x^{y^\prime}=x^{y+y^\prime}$} and $x^0$ is the unit element. We call $x^y$ the quasi-monomial with quasi-exponent $y\in E$. The Weyl group $W_0$ acts on $\mathbf{F}[E]$
by $\mathbf{F}$-algebra automorphisms by
\[
w(x^y):=x^{wy}
\]
for $w\in W_0$ and $y\in E$.

For any subset $Z\subseteq E$, we write
\[
\mathbf{F}[Z]:=\bigoplus_{y\in Z}\mathbf{F}x^y
\]
for the subspace of $\mathbf{F}[E]$ spanned by the quasi-monomials $x^y$, $y\in Z$.

The subspace $\mathbf{F}[P^\vee]$ is a $W_0$-stable subalgebra of $\mathbf{F}[E]$, which we call the subalgebra of Laurent polynomials in $\mathbf{F}[E]$. For $\mu\in P^\vee$, we say that $x^\mu$ is a monomial with exponent $\mu\in P^\vee$.
For generic multiplicity parameter $\mathbf{k}\in\mathcal{K}^{{\rm res}}$, the nonsymmetric Macdonald polynomials form a~basis of $\mathbf{F}[P^\vee]$.
In this paper, we are focussing on the theory for the extended set $\mathcal{K}$ of multiplicity parameters, in which case the corresponding nonsymmetric Macdonald--Koornwinder polynomials form a basis of the $W_0$-stable subalgebra $\mathbf{F}\smash{\bigl[Q^\vee\bigr]}$ of $\mathbf{F}[P^\vee]$.

Let $E/W$ be the set of $W$-orbits in $E$. For a $W$-orbit $\mathcal{O}\in E/W$, we denote by $c^{\mathcal{O}}$ the unique point in the intersection of $\mathcal{O}$
and the closure $\overline{C_+}$ of the fundamental alcove. Note that the $W$-orbit in $E$ containing the origin is $Q^\vee$.
A~$W$-orbit $\mathcal{O}$ in $E$ has a finite number of $\tau\smash{\bigl(Q^\vee\bigr)}$-orbits $\tau\smash{\bigl(Q^\vee\bigr)}y_i$, $1\leq i\leq N$, hence the corresponding space $\mathbf{F}[\mathcal{O}]$ of quasi-polynomials with quasi-exponents in $\mathcal{O}$ is a free $\mathbf{F}\smash{\bigl[Q^\vee\bigr]}$-module of finite rank,
\begin{equation}\label{FOpoldec}
\mathbf{F}[\mathcal{O}]=\bigoplus_{i=1}^N\mathbf{F}\smash{\bigl[Q^\vee\bigr]}x^{y_i}.
\end{equation}
Furthermore, the
space $\mathbf{F}[E]$ of quasi-polynomials decomposes as
\begin{equation}\label{Porbit}
\mathbf{F}[E]=\bigoplus_{\mathcal{O}\in E/W}\mathbf{F}[\mathcal{O}].
\end{equation}

For generic multiplicity functions $\mathbf{k}\in\mathcal{K}^{{\rm res}}$ and an arbitrary $W$-orbit $\mathcal{O}$, quasi-polynomial extensions of the nonsymmetric Macdonald polynomials were introduced in~\cite{SSV2}. They depend on a generic dilation parameter $q_\varphi\in\mathbf{F}^\times$ and additional $\mathcal{O}$-dependent representation parameters, and they form a basis of $\mathbf{F}[\mathcal{O}]$. For $\mathcal{O}=Q^\vee$, they are the nonsymmetric Macdonald polynomials.
The goal of this paper is to extend this result to multiplicity functions $\mathbf{k}$ in $\mathcal{K}$.

This boils down to introducing the Koornwinder $\bigl({\rm C}^\vee{\rm C}_r\bigr)$ analogs of the quasi-polynomial extensions of the type ${\rm C}_r$ nonsymmetric Macdonald polynomials from~\cite{SSV2}. They will now depend on five (four in case of $r=1$) multiplicity parameters instead of two (one in case of~${r=1}$). To stay close to the notations and results from~\cite{SSV2}, we will give a uniform treatment
of the theory for multiplicity functions $\mathbf{k}\in\mathcal{K}$ when the root system $\Phi_0$ is of arbitrary type. For type ${\rm C}_r$, the results are made more concrete in Sections~\ref{1.2}--\ref{1.4}.

\subsection{The affine Hecke algebra}\label{AHAsection}

The affine Hecke algebra $H=H(\mathbf{k})$ is the unital associative $\mathbf{F}$-algebra with generators $T_j$, $0\leq j\leq r$, and relations
\begin{enumerate}\itemsep=0pt\samepage
\item[(a)] The $(W,\{s_0,\dots,s_r\})$-braid relations for $T_0,\dots,T_r$,
\item[(b)] $(T_j-k_j)\bigl(T_j+k_j^{-1}\bigr)=0$ for $j=0,\dots,r$.
\end{enumerate}\pagebreak

\noindent
Here (a) means{\samepage
\[
T_iT_jT_i\cdots=T_jT_iT_j\cdots,\qquad 0\leq i\not=j\leq r,
\]
with on each side $m_{ij}$ terms, where $m_{ij}$ is the order of $s_is_j$ in $W$.}

A reduced expression of $g\in W$ is an expression $g=s_{j_1}\cdots s_{j_{\ell(g)}}$ of $g$ as product of simple reflections with $\ell(g)$ minimal.
The length function $\ell\colon W\rightarrow\mathbb{Z}_{\geq 0}$ satisfies $\ell(g)=\#\Pi(g)$, with
\[
\Pi(g):=\Phi^+\cap g^{-1}\Phi^-.
\]
The braid relations ensure that the element
\[
T_g:=T_{j_1}\cdots T_{j_{\ell(g)}}
\]
in $H$ does not depend on the choice of reduced expression $g=s_{j_1}\cdots s_{j_{\ell(g)}}$, and $\{T_g\}_{g\in W}$ is a~basis of $H$.

The finite Hecke algebra $H_0=H_0(\mathbf{k})$ is the subalgebra of $H$ generated by $T_1,\dots,T_r$. The defining relations of $H_0$ in terms of $T_1,\dots,T_r$
are the $(W_0,\{s_1,\dots,s_r\})$-braid relations and the quadratic relations $(T_i-k_i)\bigl(T_i+k_i^{-1}\bigr)=0$ for $i=1,\dots,r$.
For $w\in W_0\subset W$, a~reduced expression $w=s_{i_1}\cdots s_{i_\ell}$ in $W$ can be chosen with simple reflections from $W_0$, i.e., with $1\leq i_j\leq r$. Furthermore,
$\Pi(w)=\Phi_0^+\cap w^{-1}\Phi_0^-$ for $w\in W_0$, and $\{T_w\}_{w\in W_0}$ is a basis of $H_0$.

Define $\chi\colon \Phi_0\rightarrow \{\pm 1\}$ by
\begin{equation}\label{chi}
\chi(\alpha):=
\begin{cases}
\hphantom{-}1&\text{if}\ \alpha\in\Phi_0^{+},\\
-1&\text{if}\ \alpha\in\Phi_0^-.
\end{cases}
\end{equation}
In other words, $\chi=\chi_+-\chi_-$ with
\begin{equation*}%\label{chipm}
\chi_{\pm}:=\chi_{\Phi_0^{\pm}}\colon\ \Phi_0\rightarrow \{0,1\}
\end{equation*}
the indicator function of $\Phi_0^{\pm}$ in $\Phi_0$.
For $j=0,\dots,r$ and $g\in W$, we have
\begin{equation*}%\label{lengthformula}
\ell(s_jg)=\ell(g)+\chi\bigl(g^{-1}\alpha_j\bigr),
\end{equation*}
and hence
\begin{equation}\label{TjTw}
T_jT_g=\chi_-\bigl(g^{-1}\alpha_j\bigr)\bigl(k_j-k_j^{-1}\bigr)T_g+T_{s_jg}
\end{equation}
in the affine Hecke algebra $H$.

We now describe the Bernstein decomposition of $H$ (see~\cite{Lu} for details).
Let $H^\times$ be the group of units in $H$. There exists a unique group homomorphism
\[
Q^\vee\rightarrow H^\times,\qquad \lambda\mapsto Y^\lambda
\]
such that $Y^\lambda=T_{\tau(\lambda)}$ for \smash{$\lambda\in\overline{E_+}\cap Q^\vee$}.
The resulting algebra map
\begin{equation}\label{injY}
\mathbf{F}\smash{\bigl[Q^\vee\bigr]}\hookrightarrow H,\qquad p\mapsto p(Y),
\end{equation}
mapping $x^\lambda$ to $Y^\lambda$ for all $\lambda\in Q^\vee$, is injective. The image of the embedding is denoted by~\smash{$\mathbf{F}_Y\smash{\bigl[Q^\vee\bigr]}$}.
The multiplication map of $H$ restricts to a linear isomorphism
\[
H_0\otimes\mathbf{F}_Y\smash{\bigl[Q^\vee\bigr]}\overset{\sim}{\longrightarrow} H.
\]
The commutation relations between elements in $H_0$ and \smash{$\mathbf{F}_Y\smash{\bigl[Q^\vee\bigr]}$} are described (and determined) by the {Bernstein--Lusztig} cross relations
\begin{equation}\label{crossY}
Y^\lambda T_i-T_i Y^{s_i\lambda}=\Biggl(\frac{\widetilde{\mathbf{k}}_{\alpha_i}-\widetilde{\mathbf{k}}_{\alpha_i}^{-1}+
\Bigl(\widetilde{\mathbf{k}}_{\frac{\alpha_i}{2}}-\widetilde{\mathbf{k}}_{\frac{\alpha_i}{2}}^{-1}\Bigr)Y^{-\alpha_i^\vee}}{1-Y^{-2\alpha_i^\vee}}\Biggr)\bigl(Y^\lambda-Y^{s_i\lambda}\bigr)
\end{equation}
for $i=1,\dots,r$ and $\lambda\in Q^\vee$.

The right-hand side of~\eqref{crossY} appears to be in the quotient field $\mathbf{F}_Y\smash{\bigl(Q^\vee\bigr)}$ of $\mathbf{F}_Y\smash{\bigl[Q^\vee\bigr]}$, but it lies
in $\mathbf{F}_Y\smash{\bigl[Q^\vee\bigr]}$. Indeed, if $\alpha_i\smash{\bigl(Q^\vee\bigr)}=\mathbb{Z}$, then
\smash{$\widetilde{\mathbf{k}}_{\frac{\alpha_i}{2}}=\widetilde{\mathbf{k}}_{\alpha_i}=k_i$}, hence the right-hand side of~\eqref{crossY} reduces to
\[
\bigl(k_i-k_i^{-1}\bigr)\biggl(\frac{Y^\lambda-Y^{s_i\lambda}}{1-Y^{-\alpha_i^\vee}}\biggr)=
\bigl(k_i-k_i^{-1}\bigr)Y^\lambda\biggl(\frac{1-Y^{-\alpha_i(\lambda)\alpha^\vee}}{1-Y^{-\alpha_i^\vee}}\biggr),
\]
which lies in $\mathbf{F}_Y\smash{\bigl[Q^\vee\bigr]}$ since $\alpha_i(\lambda)\in\mathbb{Z}$. If $\alpha_i\smash{\bigl(Q^\vee\bigr)}=\mathbb{Z}_e$, then $\Phi_0$ is of type ${\rm C}_r$, $r\geq 1$,
and $i=r$, in view of our convention that $\alpha_r$ is a long root.
The right-hand side of~\eqref{crossY} then reads
\[
\bigl(k_r-k_r^{-1}+\bigl(k_0-k_0^{-1}\bigr)Y^{-\alpha_r^\vee}\bigr)Y^\lambda\biggl(\frac{1-Y^{-\alpha_r(\lambda)\alpha_r^\vee}}{1-Y^{-2\alpha_r^\vee}}\biggr),
\]
which lies in $\mathbf{F}_Y\smash{\bigl[Q^\vee\bigr]}$ since $\alpha_r(\lambda)\in \mathbb{Z}_e$.

Denote by
\[
\mathbf{F}\smash{\bigl[Q^\vee\bigr]}^{W_0}\subset\mathbf{F}\smash{\bigl[Q^\vee\bigr]}
\]
the subalgebra of $W_0$-invariant elements in $\mathbf{F}\smash{\bigl[Q^\vee\bigr]}$, and \smash{$\mathbf{F}_Y\bigl[Q^\vee\bigr]^{W_0}$} for its image in $H$ under the embedding~\eqref{injY}. Then
\[
Z(H)=\mathbf{F}_Y\smash{\bigl[Q^\vee\bigr]}^{W_0},
\]
where $Z(H)$ denotes the center of $H$.

%%%%%%%%%%%%%%%%%%%%%%%%%%%%%%%%%%%%%%%%%%%
\subsection{The double affine Hecke algebra}\label{2.5}
%%%%%%%%%%%%%%%%%%%%%%%%%%%%%%%%%%%%%%%%%%%
Fix a parameter $q_\varphi\in\mathbf{F}^\times$ and set
\[
q_\alpha:=q_\varphi^{\|\varphi\|^2/\|\alpha\|^2}\qquad \text{for}\ \alpha\in\Phi_0.
\]
It equals either $q_\varphi$, $q_\varphi^2$ or $q_\varphi^3$ (see~\cite[Section~9.4, Table~1]{Hu}), and $q_{w\alpha}=q_\alpha$ for all $\alpha\in\Phi_0$.
We like to think of $q_\varphi$ as being equal to \smash{$q^{2/\|\varphi\|^2}$} for $q$ in some field extension of $\mathbf{F}$ (this is done in~Sections~\ref{1.2}--\ref{1.4}, where we described the results explicitly for $\Phi_0$ of type ${\rm C}_r$).
To circumvent field extensions, we will introduce instead a group homomorphism $q\colon \smash{\frac{2}{\|\varphi\|^2}}\mathbb{Z}\rightarrow\mathbf{F}^\times$, $m\mapsto q^m$, with
\begin{equation*}%\label{qpower}
q^{m}:=q_\varphi^{m\|\varphi\|^2/2}\qquad \text{for}\ m\in\tfrac{2}{\|\varphi\|^2}\mathbb{Z}.
\end{equation*}
Note that $q^m\in\mathbf{F}$ is meaningful for $m\in \bigl\langle Q^\vee, Q^\vee\bigr\rangle$ since
$\bigl\langle Q^\vee, Q^\vee\bigr\rangle\subseteq\frac{2}{\|\varphi\|^2}\mathbb{Z}$.

%%%%%%%%%%%%%%%%%%%%%%%%%
\begin{Definition}\label{T_def}
We denote by $\mathbf{T}$ the $\mathbf{F}$-torus of rank $r$ consisting of the group homomorphisms~$Q^\vee\rightarrow\mathbf{F}^\times$.
\end{Definition}
%%%%%%%%%%%%%%%%%%%%%%%%
The abelian group structure on $\mathbf{T}$ is by pointwise multiplication,
\[ (st)^\mu:=s^\mu t^\mu,\qquad s,t\in\mathbf{T},\ \mu\in Q^\vee,
\]
where $t^\mu\in\mathbf{F}^\times$ denotes the value of $t\in\mathbf{T}$ at $\mu\in Q^\vee$.
%%%%%%%%%%%%%%%%%%%%%%%%%%%%
\begin{Remark}
Consider the root system $\Phi_0=\{\pm(\epsilon_i\pm\epsilon_j)\}_{1\leq i<j\leq r}\cup\{\pm 2\epsilon_i\}_{i=1}^r$ of type ${\rm C}_r$, with~$\{\epsilon_1,\dots,\epsilon_r\}$ the standard orthonormal basis of $\mathbb{R}^r$. Its co-root lattice $Q^\vee$
equals $\mathbb{Z}^r$. In the introduction, we identified the corresponding $\mathbf{F}$-torus $\mathbf{T}$ with $(\mathbf{F}^\times)^r$ by the isomorphism
\[
\mathbf{T}\overset{\sim}{\longrightarrow}(\mathbf{F}^\times)^r,\qquad t\mapsto (t^{\epsilon_1},\dots,t^{\epsilon_r}),
\]
cf.~\eqref{F_torus_CC} and~\eqref{F_torus_CCC}.
\end{Remark}


For $\lambda\in Q^\vee$, we define the torus element
\[
q^\lambda\in\mathbf{T}
\]
by $\mu\mapsto q^{\langle \lambda,\mu\rangle}$, $\mu\in Q^\vee$. Then $\mathbf{T}$ admits a left $W$-action by
\begin{align}
&(wt)^\mu:=t^{w^{-1}\mu},\nonumber\\
&(\tau(\lambda)t)^\mu:=(q^\lambda t)^\mu=q^{\langle\lambda,\mu\rangle}t^\mu\label{Waction}
\end{align}
for $t\in\mathbf{T}$, $w\in W_0$ and $\lambda,\mu\in Q^\vee$.

We will view a polynomial $p=\sum_\mu d_\mu x^\mu\in\mathbf{F}\smash{\bigl[Q^\vee\bigr]}$ as regular function on $\mathbf{T}$ by
\[
p(t):=\sum_\mu d_\mu t^\mu\qquad \text{for}\ t\in\mathbf{T}.
\]
The formula
\[
(gp)(t):=p\bigl(g^{-1}t\bigr)
\]
for $p\in\mathbf{F}\smash{\bigl[Q^\vee\bigr]}$, $g\in W$ and $t\in\mathbf{T}$ then turns $\mathbf{F}\smash{\bigl[Q^\vee\bigr]}$ into a $W$-module, with $W$ acting by algebra automorphisms.
Concretely, the action on the basis of monomials is given by
\begin{equation}\label{Wa}
w(x^\mu)=x^{w\mu},\qquad
\tau(\lambda)(x^\mu)=q^{-\langle\mu,\lambda\rangle}x^\mu
\end{equation}
for $w\in W_0$ and $\lambda,\mu\in Q^\vee$.

Note that by~\eqref{ssplit},
\begin{equation}\label{ssplitaction}
s_a(x^\mu)=q_{\overline{a}}^{-a(0)\overline{a}(\mu)}x^{s_{\overline{a}}\mu}
\end{equation}
for $a\in\Phi$ and $\mu\in Q^\vee$. In particular,
\begin{equation*}%\label{s0actionuntwisted}
s_0(x^\mu)=q_\varphi^{\varphi(\mu)}x^{s_\varphi\mu}.
\end{equation*}

In various computations, it is convenient to use co-roots of affine roots and incorporate $q$-powers in the exponents of the monomials. The co-root $b^\vee$ of $b\in\Phi$ is defined by
\[
b^\vee:=\biggl(\overline{b}^\vee,\frac{2b(0)}{\|\overline{b}\|^2}\biggr)\in E\times\mathbb{R}.
\]
The resulting set $\Phi^\vee:=\{b^\vee\}_{b\in\Phi}$ of co-roots is an affine root system in $E\times\mathbb{R}$~\cite{MacAff}. Note that
\begin{equation}\label{dualactionW}
(s_ab)^\vee=b^\vee-\overline{a}\bigl(\overline{b}^\vee\bigr)a^\vee\qquad \text{for}\, a,b\in\Phi
\end{equation}
in view of
\eqref{WPhiaction} and the fact that \smash{$\frac{2}{\|\beta\|^2}\beta(\alpha^\vee)=\frac{2}{\|\alpha\|^2}\alpha(\beta^\vee)$} for $\alpha,\beta\in\Phi_0$ (indeed, both sides are equal to $\langle\alpha^\vee,\beta^\vee\rangle$ by~\eqref{corootdef}).

We set
\[
x^{\widehat{\mu}}:=q^mx^\mu\in\mathbf{F}\smash{\bigl[Q^\vee\bigr]}\qquad\textup{for}\ \widehat{\mu}=(\mu,m)\in Q^\vee\times \tfrac{2}{\|\varphi\|^2}\mathbb{Z}
\]
and we write $t^{\widehat{\mu}}$ for the evaluation of $x^{\widehat{\mu}}\in\mathbf{F}\smash{\bigl[Q^\vee\bigr]}$ at $t\in\mathbf{T}$ \big(so in particular, $t^{\widehat{\mu}}=q^mt^\mu$\big).
Note that \smash{$\Phi^\vee\subset Q^\vee\times \frac{2}{\|\varphi\|^2}\mathbb{Z}$}, hence \smash{$x^{b^\vee}$} makes sense for all $b\in\Phi$. Concretely,
\begin{equation}\label{arootreduction}
x^{b^\vee}:=q_{\overline{b}}^{b(0)}x^{\overline{b}^\vee}\in\mathbf{F}\smash{\bigl[Q^\vee\bigr]},
\end{equation}
which reduces to the monomial \smash{$x^{\beta^\vee}$} in case $b=(\beta,0)$.

Formula~\eqref{ssplitaction} can then be rewritten as
\begin{equation}\label{ssplitaction2}
s_a(x^\mu)=x^{\mu-\overline{a}(\mu)a^\vee}.
\end{equation}
We furthermore have
\begin{equation}\label{awcomp}
g(x^{b^\vee})=x^{(gb)^\vee}\qquad \textup{for}\ g\in W\ \text{and}\ b\in\Phi.
\end{equation}
Indeed, it suffices to check~\eqref{awcomp} for $w=s_a$, $a\in\Phi$. By~\eqref{arootreduction} and~\eqref{ssplitaction2}, we have
\[
s_a\bigl(x^{b^\vee}\bigr)=x^{b^\vee-\overline{a}(\overline{b}^\vee)a^\vee}
\]
which equals \smash{$x^{(s_ab)^\vee}$} by~\eqref{dualactionW}.

\begin{Definition}[\cite{ChKZ2,Sa}]
The double affine Hecke algebra $\mathbb{H}=\mathbb{H}(\mathbf{k};q_\varphi)$ is the unital associative $\mathbf{F}$-algebra generated by
$T_j$, $0\leq j\leq r$, and $x^\mu$, $\mu\in Q^\vee$, subject to the following relations:
\begin{enumerate}\itemsep=0pt
\item[(a)] the $(W,\{s_0,\dots,s_r\})$-braid relations for $T_0,\dots,T_r$,
\item[(b)] the quadratic relations $(T_j-k_j)\bigl(T_j+k_j^{-1}\bigr)=0$ for $j=0,\dots,r$,
\item[(c)] $x^\mu x^\nu=x^{\mu+\nu}$, $\mu,\nu\in Q^\vee$, and $x^0$ is the unit element of $\mathbb{H}$,
\item[(d)] the cross relations
\begin{equation}\label{crossX}
T_jx^\mu-s_j(x^\mu)T_j=\Biggl(\frac{k_j-k_j^{-1}+\bigl(u_j-u_j^{-1}\bigr)x^{\alpha_j^\vee}}{1-x^{2\alpha_j^\vee}}\Biggr)(x^\mu-s_j(x^\mu))
\end{equation}
for $j=0,\dots,r$ and $\mu\in Q^\vee$.
\end{enumerate}
\end{Definition}
%%%%%%%%%%%%%%%%%%%%%%%%%%%%%%%%%%%%%%%%%%%%%
The Poincar{\'e}--Birkhoff--Witt (PBW) theorem for $\mathbb{H}$ states that the canonical algebra maps $H\rightarrow \mathbb{H}$ and $\mathbf{F}\smash{\bigl[Q^\vee\bigr]}\rightarrow
\mathbb{H}$
are embeddings, and that
the multiplication map of $\mathbb{H}$ restricts to a~linear isomorphism
\[
\mathbf{F}\smash{\bigl[Q^\vee\bigr]}\otimes H\overset{\sim}{\longrightarrow} \mathbb{H}.
\]
By the Bernstein presentation of the affine Hecke algebra (see Section~\ref{AHAsection}), the subalgebra
\[
H^X\subset\mathbb{H}
\]
generated by $\mathbf{F}\smash{\bigl[Q^\vee\bigr]}$ and $H_0$ is isomorphic to the affine Hecke algebra $\widetilde{H}:=H(\widetilde{\mathbf{k}})$.

The double affine Hecke algebra with dual multiplicity parameters will be denoted by
\[
\widetilde{\mathbb{H}}:=\mathbb{H}\bigl(\widetilde{\mathbf{k}},q_\varphi\bigr).
\]
To keep the notations manageable, we will use the same notations $T_j$, $T_g$, $Y^\mu$, $x^\mu$ in both $\mathbb{H}$ and~$\widetilde{\mathbb{H}}$.

The duality anti-isomorphism~\cite{ChBook,Sa} is the unique anti-algebra isomorphism
\[
\delta=\delta_{\mathbf{k}}\colon\ \mathbb{H}\overset{\sim}{\longrightarrow}\widetilde{\mathbb{H}}
\]
satisfying
\[
\delta(T_i)=T_i,\qquad \delta\bigl(Y^\lambda\bigr)=x^{-\lambda},\qquad \delta(x^\mu)=Y^{-\mu}
\]
for $i=1,\dots,r$ and $\lambda,\mu\in Q^\vee$. Its inverse is $\widetilde{\delta}:=\delta_{\widetilde{\mathbf{k}}}$.

\section{Quasi-polynomial representations}\label{qpsect}

\subsection[The quasi-polynomial representation of H\^{}X]{The quasi-polynomial representation of $\boldsymbol{H^X}$}\label{qpHX}

In this subsection, we introduce the quasi-polynomial representation of the dual affine Hecke algebra $H^X$.
In case $\mathbf{k}\in\mathcal{K}^{{\rm res}}$,
the representation we obtain was derived before in~\cite{SSV2, SSV1}.

Let
\[
\lfloor\cdot\rfloor\colon \ \mathbb{R}\rightarrow \mathbb{Z}
\]
be the floor function, so $\lfloor s\rfloor$ is the largest integer $\leq s$. We denote by
\[
\lfloor\cdot\rfloor_e\colon \ \mathbb{R}\rightarrow\mathbb{Z}_e,\qquad
\lfloor\cdot\rfloor_o\colon \ \mathbb{R}\rightarrow\mathbb{Z}_o
\]
the functions which map $s\in\mathbb{R}$ to
the smallest even and odd integer $\leq s$, respectively. Note~that
\[
\lfloor s\rfloor_e:=2\lfloor s/2\rfloor,\qquad
\lfloor s\rfloor_o=\lfloor s+1\rfloor_e-1
\]
for $s\in\mathbb{R}$.

\begin{Definition}\label{truncation}
For $a\in\Phi$, define the even and odd truncated divided difference operator $\nabla_a=\nabla_a(\mathbf{k})\in\textup{End}(\mathbf{F}[E])$ by
\begin{align}
&\nabla_a^e(x^y):=\Biggl(\frac{1-x^{-\lfloor\overline{a}(y)\rfloor_ea^\vee}}{1-x^{2a^\vee}}\Biggr)x^y,\nonumber\\
&\nabla_a^o(x^y):=\Biggl(\frac{x^{a^\vee}-x^{-\lfloor\overline{a}(y)\rfloor_oa^\vee}}{1-x^{2a^\vee}}\Biggr)x^y\label{Ai}
\end{align}
for $y\in E$. We furthermore write $\nabla_j^e:=\nabla_{\alpha_j}^e$ and $\nabla_j^o:=\nabla_{\alpha_j}^o$ for $j=0,\dots,r$.
\end{Definition}

Truncation in Definition~\ref{truncation} refers to the fact that the real numbers $\overline{a}(y)$ in formula~\eqref{Ai} are truncated using the even and odd floor operations.
These truncations are necessary to turn the two quotients in~\eqref{Ai} into well defined elements in $\mathbf{F}\smash{\bigl[Q^\vee\bigr]}$. Note that $\nabla_a^e$ and $\nabla_a^o$ depend on~$q_\varphi$ when $a\in\Phi\setminus\Phi_0$, in view of~\eqref{arootreduction}. In this subsection, we only need the truncated divided difference operators for $a\in\Phi_0$, and there will be no dependence on $q_\varphi$.

We write
\begin{equation}\label{Aisum}
\nabla_a:=\nabla_a^e+\nabla_a^o
\end{equation}
for the sum of the even and odd truncated divided difference operator, and $\nabla_j:=\nabla_{\alpha_j}$ for $j=0,\dots,r$. Then
\begin{equation}\label{Aires}
\nabla_a(x^y)=\Biggl(\frac{1-x^{-\lfloor\overline{a}(y)\rfloor a^\vee}}{1-x^{a^\vee}}\Biggr)x^y
\end{equation}
for $y\in E$ since
\begin{equation*}%\label{2seteq}
\{\lfloor s\rfloor_e,\lfloor s\rfloor_o\}=\{\lfloor s\rfloor,\lfloor s\rfloor-1\}
\end{equation*}
as unordered $2$-sets for any $s\in\mathbb{R}$.
The truncated difference operator $\nabla_a\in\textup{End}(\mathbf{F}[E])$ was introduced before in~\cite[Section~4.2]{SSV2}.

The link of the various truncated divided difference operators to the usual divided difference operator is as follows (the first part of the lemma was observed before in~\cite[Lemma~4.4]{SSV2}).
%%%%%%%%%%%%%%%%%%%%%%%%%%%%%
\begin{Lemma}\label{polredNabla}
\hfill
\begin{enumerate}\itemsep=0pt
\item[$(1)$] If $a\in\Phi$, then
\[
\nabla_a(x^\mu)=\frac{x^\mu-s_a(x^{\mu})}{1-x^{a^\vee}}
\]
for $\mu\in Q^\vee$.
\item[$(2)$] If $\Phi_0$ is of type ${\rm C}_r$, $r\geq 1$, and if $a\in\Phi$ is an affine root such that $\overline{a}\in\Phi_0^\ell$ \textup{(}in other words,
$a\in W\alpha_0\sqcup W\alpha_r$\textup{)},
then
\[
\nabla_a^e(x^\mu)=\frac{x^\mu-s_a(x^{\mu})}{1-x^{2a^\vee}}=x^{-a^\vee}\nabla_a^o(x^\mu)
\]
for $\mu\in Q^\vee$.
\end{enumerate}
\end{Lemma}
%%%%%%%%%%%%%%%%%%%%%%%%%%%%
\begin{proof}
(1) This is immediate from~\eqref{ssplitaction2} and the fact that $\overline{a}\smash{\bigl(Q^\vee\bigr)}\subseteq\mathbb{Z}$.\\
(2)
Under these assumptions, we have $\overline{a}\smash{\bigl(Q^\vee\bigr)}=\mathbb{Z}_e$, hence
\[
\lfloor \overline{a}(\mu)\rfloor_e=\overline{a}(\mu)=\lfloor \overline{a}(\mu)\rfloor_o+1
\]
and the result follows again from~\eqref{ssplitaction2}.
\end{proof}

We will often make use of the indicator functions $\chi_{\mathbb{Z}_e}, \chi_{\mathbb{Z}_o}\colon \mathbb{R}\rightarrow\{0,1\}$ of the
subsets of even and odd integers, respectively. We will denote them by $\chi_e$ and $\chi_o$, respectively.

\begin{Theorem}\label{theoremH}
The formulas
\begin{align}
&\pi(x^\mu)x^y:=x^{y+\mu},\nonumber\\
&\pi(T_i)x^y:=k_i^{\chi_e(\alpha_i(y))}u_i^{\chi_o(\alpha_i(y))}x^{s_iy}+\bigl(k_i-k_i^{-1}\bigr)\nabla_i^e(x^y)+\bigl(u_i-u_i^{-1}\bigr)\nabla_i^o(x^y)\label{qprepHX}
\end{align}
for $\mu\in Q^\vee$, $i\in\{1,\dots,r\}$ and $y\in E$
define a representation $\pi\colon H^X\rightarrow\textup{End}(\mathbf{F}[E])$.
\end{Theorem}
%%%%%%%%%%%%%%%%%%%%%%%%%%%%%%%%%%%%%%%%%%%%%%%%%%%%%%%%%
\begin{proof}
Formula~\eqref{qprepHX} uniquely defines linear operators $\pi(x^\mu)$ and $\pi(T_i)$ on $\mathbf{F}[E]$, which in turn restrict to linear operators
\begin{equation}\label{piO}
\pi^{\mathcal{O}}(x^\mu):=\pi(x^\mu)\vert_{\mathbf{F}[\mathcal{O}]},\qquad \pi^{\mathcal{O}}(T_i):=\pi(T_i)\vert_{\mathbf{F}[\mathcal{O}]}
\end{equation}
on $\mathbf{F}[\mathcal{O}]$ for every $W$-orbit $\mathcal{O}$ in $E$. In view of~\eqref{Porbit} it thus suffices to show that the linear operators~\eqref{piO} define a representation $\pi^{\mathcal{O}}=\pi(\cdot)\vert_{\mathbf{F}[\mathcal{O}]}\colon H^X\rightarrow\textup{End}(\mathbf{F}[\mathcal{O}])$.

Consider $H^X$ as left regular $H^X$-module. By~\eqref{TjTw} and~\eqref{crossX} the $H^X$-action can be written down explicitly relative to the basis
\[
\bigl\{x^\lambda T_w\mid \lambda\in Q^\vee,\, w\in W_0\bigr\}
\]
of $H^X$. The resulting formulas are
\begin{gather}
x^\mu x^\lambda T_w=x^{\lambda+\mu}T_w,\nonumber\\
T_ix^\lambda T_w=x^{s_i\lambda}T_{s_iw}+\bigl(k_i-k_i^{-1}\bigr)\Biggl(\frac{1-x^{(2\chi_-\bigl(w^{-1}\alpha_i\bigr)-\alpha_i(\lambda))\alpha_i^\vee}}{1-x^{2\alpha_i^\vee}}\Biggr)x^\lambda T_w\nonumber\\
\hphantom{T_ix^\lambda T_w=}{}
+(u_i-u_i^{-1})\Biggl(\frac{x^{\alpha_i^\vee}-x^{(1-\alpha_i(\lambda))\alpha_i^\vee}}{1-x^{2\alpha_i^\vee}}\Biggr)x^{\lambda}T_w
\label{regularHX}
\end{gather}
for $\lambda,\mu\in Q^\vee$, $w\in W_0$ and $i=1,\dots,r$.

Let $\kappa_w^\mathcal{O}\in\mathbf{F}^\times$, $w\in W_0$, be a collection of nonzero scalars.
Consider the surjective linear map
\begin{equation}\label{psicmap}
\psi^\mathcal{O}\colon\ H^X\twoheadrightarrow \mathbf{F}[\mathcal{O}],\qquad \psi^\mathcal{O}\bigl(x^\lambda T_w\bigr):=\kappa_w^\mathcal{O}x^{\lambda+wc^{\mathcal{O}}},\quad \lambda\in Q^\vee,\ w\in W_0.
\end{equation}
We will fine tune the scalars $\kappa_w^\mathcal{O}\in\mathbf{F}^\times$, $w\in W_0$, in such a way that the kernel of $\psi^\mathcal{O}$
is a left ideal in $H^X$. This will allow us to push the left regular $H^X$-action through $\psi^\mathcal{O}$, giving rise to a
$H^X$-action on $\mathbf{F}[\mathcal{O}]$. We then show that the resulting action of $x^\mu$ and $T_i$ on $\mathbf{F}[\mathcal{O}]$ is by the linear operators $\pi^\mathcal{O}(x^\mu)$ and $\pi^\mathcal{O}(T_i)$, which completes the proof of the theorem.

For any choice of scalars $\kappa_w^\mathcal{O}\in\mathbf{F}^\times$, $w\in W_0$, the kernel of $\psi^\mathcal{O}$ is invariant under left multiplication by $\mathbf{F}\smash{\bigl[Q^\vee\bigr]}$.
Note that the kernel $\psi^\mathcal{O}$ is invariant under left multiplication by $T_i$ if there exists a linear operator $\mathcal{D}_i^\mathcal{O}\in\textup{End}(\mathbf{F}[\mathcal{O}])$ such that
\begin{equation*}%\label{DiIntro}
\mathcal{D}_i^\mathcal{O}\bigl(\psi^\mathcal{O}(h)\bigr)=\psi^\mathcal{O}(T_ih)
\end{equation*}
for all $h\in H^X$.
By~\eqref{regularHX}, we have
\begin{align}
\psi^\mathcal{O}\bigl(T_ix^\lambda T_w\bigr)={}&\kappa_{s_iw}^\mathcal{O}s_i\bigl(x^{\lambda+wc^{\mathcal{O}}}\bigr)
+ \kappa_w^\mathcal{O}\bigl(k_i-k_i^{-1}\bigr)\Biggl(\frac{1-x^{(2\chi_-\bigl(w^{-1}\alpha_i\bigr)-\alpha_i(\lambda))\alpha_i^\vee}}{1-x^{2\alpha_i^\vee}}\Biggr)x^{\lambda+wc^{\mathcal{O}}}\nonumber\\
&{}+
\kappa_w^\mathcal{O}\bigl(u_i-u_i^{-1}\bigr)\Biggl(\frac{x^{\alpha_i^\vee}-x^{(1-\alpha_i(\lambda))\alpha_i^\vee}}{1-x^{2\alpha_i^\vee}}\Biggr)x^{\lambda+wc^{\mathcal{O}}},
\label{spoint}
\end{align}
so we look for conditions on the scalars $\kappa_w^\mathcal{O}$ such that~\eqref{spoint} can be expressed as
a linear operator $\mathcal{D}_i^\mathcal{O}\in\textup{End}(\mathbf{F}[\mathcal{O}])$ acting on $\psi^\mathcal{O}\bigl(x^\lambda T_w\bigr)=\kappa_w^\mathcal{O}x^{\lambda+wc^\mathcal{O}}$ for all $\lambda\in Q^\vee$ and $w\in W_0$.

To simplify notations, we write
\[
y=\lambda+wc^\mathcal{O}
\]
with $w\in W_0$ and $\lambda\in Q^\vee$.
We first prove that
\begin{gather}
\frac{1}{\kappa_w^\mathcal{O}} \psi^\mathcal{O}\bigl(T_ix^\lambda T_w\bigr)
=\biggl(\frac{\kappa_{s_iw}^\mathcal{O}}{\kappa_w^\mathcal{O}}
+\bigl(k_i-k_i^{-1}\bigr)\chi_-\bigl(w^{-1}\alpha_i\bigr)\chi_e\bigl(\alpha_i\bigl(wc^\mathcal{O}\bigr)\bigr)\nonumber\\
\hphantom{\frac{1}{\kappa_w^\mathcal{O}} \psi^\mathcal{O}\bigl(T_ix^\lambda T_w\bigr)=\biggl(}{}
 +\bigl(u_i-u_i^{-1}\bigr)\chi_+\bigl(w^{-1}\alpha_i\bigr)\chi_o(\alpha_i\bigl(wc^\mathcal{O}\bigr))\biggr)x^{s_iy}\nonumber\\
 \hphantom{\frac{1}{\kappa_w^\mathcal{O}} \psi^\mathcal{O}\bigl(T_ix^\lambda T_w\bigr)=}{}
 +\bigl(k_i-k_i^{-1}\bigr)\nabla_i^e(x^y)+\bigl(u_i-u_i^{-1}\bigr)\nabla_i^o(x^y)
\label{formula}
\end{gather}
by rewriting the two quotients in~\eqref{spoint} in terms of the odd and even truncated difference operators.

We first consider the proof of~\eqref{formula} when $\alpha_i\smash{\bigl(Q^\vee\bigr)}=\mathbb{Z}_e$, i.e., when $\Phi_0$ is of type ${\rm C}_r$, $r\geq 1$, and $i=r$.
Then~\eqref{W0C} and $\alpha_r\smash{\bigl(Q^\vee\bigr)}=\mathbb{Z}_e$ imply that
\begin{equation*}
2\chi_-\bigl(w^{-1}\alpha_r\bigr)-\alpha_r(\lambda)=
\begin{cases}
2-\lfloor\alpha_r(y)\rfloor_e &\hbox{if }  w^{-1}\alpha_r\in\Phi_0^-\ \text{and}\ \alpha_r\bigl(wc^\mathcal{O}\bigr)=0,\\
-\lfloor\alpha_r(y)\rfloor_e &\textup{otherwise}.
\end{cases}
\end{equation*}
By~\eqref{W0C},
this can be reformulated as
\begin{equation*}%\label{sstep1}
2\chi_-\bigl(w^{-1}\alpha_r\bigr)-\alpha_r(\lambda)=
\begin{cases}
2-\lfloor\alpha_r(y)\rfloor_e &\text{if}\ \chi_-\bigl(w^{-1}\alpha_r\bigr)\chi_e\bigl(\alpha_r\bigl(wc^\mathcal{O}\bigr)\bigr)=1,\\
-\lfloor\alpha_r(y)\rfloor_e &\text{if}\ \chi_-\bigl(w^{-1}\alpha_r\bigr)\chi_e\bigl(\alpha_r\bigl(wc^\mathcal{O}\bigr)\bigr)=0.
\end{cases}
\end{equation*}
This allows us to rewrite the second line of~\eqref{spoint} for $\Phi_0$ of type ${\rm C}_r$ and $i=r$ in terms of the even truncated divided difference operator,
\begin{equation}\label{claim1}
\Biggl(\frac{1-x^{(2\chi_-(w^{-1}\alpha_r)-\alpha_r(\lambda))\alpha_r^\vee}}{1-x^{2\alpha_r^\vee}}\Biggr)x^{y}=
\nabla_r^e(x^y)+
\chi_-\bigl(w^{-1}\alpha_r\bigr)\chi_e\bigl(\alpha_r\bigl(wc^\mathcal{O}\bigr)\bigr)x^{s_ry}
\end{equation}
\big(in case $\chi_-\bigl(w^{-1}\alpha_r\bigr)\chi_e\bigl(\alpha_r\bigl(wc^\mathcal{O}\bigr)\bigr)=1$ use the fact that $\lfloor\alpha_r(y)\rfloor_e=\alpha_r(y)$\big).
To rewrite the third line of~\eqref{spoint} for $\Phi_0$ of type ${\rm C}_r$ and $i=r$, note that
\begin{equation*}
1-\alpha_r(\lambda)=
\begin{cases}
2-\lfloor\alpha_r(y)\rfloor_o &\text{if}\ \alpha_r\bigl(wc^\mathcal{O}\bigr)=1,\\
-\lfloor\alpha_r(y)\rfloor_o &\text{otherwise},
\end{cases}
\end{equation*}
which can be reformulated as
\begin{equation*}%\label{sstep2}
1-\alpha_r(\lambda)=
\begin{cases}
2-\lfloor\alpha_r(y)\rfloor_o &\text{if}\ \chi_+\bigl(w^{-1}\alpha_r\bigr)\chi_o\bigl(\alpha_r\bigl(wc^\mathcal{O}\bigr)\bigr)=1,\\
-\lfloor\alpha_r(y)\rfloor_o &\text{if}\ \chi_+\bigl(w^{-1}\alpha_r\bigr)\chi_o\bigl(\alpha_r\bigl(wc^\mathcal{O}\bigr)\bigr)=0
\end{cases}
\end{equation*}
in view of~\eqref{W0C}. It follows that
\begin{equation}\label{claim2}
\Biggl(\frac{x^{\alpha_r^\vee}-x^{(1-\alpha_r(\lambda))\alpha_r^\vee}}{1-x^{2\alpha_r^\vee}}\Biggr)x^y=
\nabla_r^o(x^y)+
\chi_+\bigl(w^{-1}\alpha_r\bigr)\chi_o\bigl(\alpha_r\bigl(wc^\mathcal{O}\bigr)\bigr)
x^{s_ry}
\end{equation}
(in case $\chi_+(w^{-1}\alpha_r)\chi_o(\alpha_r\bigl(wc^\mathcal{O}\bigr))=1$ use the fact that $\lfloor\alpha_r(y)\rfloor_o=\alpha_r(y)$).
Substituting~\eqref{claim1} and~\eqref{claim2} into formula~\eqref{spoint} for $i=r$, we obtain~\eqref{formula} for $\Phi_0$ of type ${\rm C}_r$ and for $i=r$.

We now prove~\eqref{formula} when $\alpha_i\smash{\bigl(Q^\vee\bigr)}=\mathbb{Z}$. Then $k_i=u_i$ and using that
\[
\bigl\{2\chi_-\bigl(w^{-1}\alpha_i\bigr)-\alpha_i(\lambda), 1-\alpha_i(\lambda)\bigr\}=\bigl\{\chi_-\bigl(w^{-1}\alpha_i\bigr)-\alpha_i(\lambda),1+\chi_-\bigl(w^{-1}\alpha_i\bigr)-\alpha_i(\lambda)\bigr\}
\]
as unordered $2$-sets,
formula~\eqref{spoint} simplifies to
\begin{equation}\label{spointred}
\psi^\mathcal{O}(T_ix^\lambda T_w)=\kappa_{s_iw}^\mathcal{O}x^{s_iy}
+\kappa_w^\mathcal{O}\bigl(k_i-k_i^{-1}\bigr)\Biggl(\frac{1-x^{(\chi_-\bigl(w^{-1}\alpha_i\bigr)-\alpha_i(\lambda))\alpha_i^\vee}}{1-x^{\alpha_i^\vee}}\Biggr)x^{y}.
\end{equation}
Using~\eqref{W0C}, we have for $w^{-1}\alpha_i\in\Phi_0^-$ that
\begin{equation*}
\chi_-\bigl(w^{-1}\alpha_i\bigr)-\alpha_i(\lambda)=
\begin{cases}
-\lfloor\alpha_i(y)\rfloor &\text{if}\ \chi_e\bigl(\alpha_i\bigl(wc^\mathcal{O}\bigr)\bigr)=0,\\
1-\lfloor\alpha_i(y)\rfloor &\text{if}\ \chi_e\bigl(\alpha_i\bigl(wc^\mathcal{O}\bigr)\bigr)=1
\end{cases}
\end{equation*}
and for $w^{-1}\alpha_i\in\Phi_0^+$ that
\begin{equation*}
\chi_-\bigl(w^{-1}\alpha_i\bigr)-\alpha_i(\lambda)=
\begin{cases}
-\lfloor\alpha_i(y)\rfloor &\text{if}\ \chi_o\bigl(\alpha_i\bigl(wc^\mathcal{O}\bigr)\bigr)=0,\\
1-\lfloor\alpha_i(y)\rfloor &\text{if}\ \chi_0\bigl(\alpha_i\bigl(wc^\mathcal{O}\bigr)\bigr)=1.
\end{cases}
\end{equation*}
This leads to the formula
\begin{gather*}
\Biggl(\frac{1-x^{(\chi_-\bigl(w^{-1}\alpha_i\bigr)-\alpha_i(\lambda))\alpha_i^\vee}}{1-x^{\alpha_i^\vee}}\Biggr)x^{y}\\
\qquad{}{}=
\nabla_i(x^y)+\bigl(\chi_-\bigl(w^{-1}\alpha_i\bigr)\chi_e\bigl(\alpha_i\bigl(wc^\mathcal{O}\bigr)\bigr)+\chi_+\bigl(w^{-1}\alpha_i\bigr)\chi_o\bigl(\alpha_i\bigl(wc^\mathcal{O}\bigr)\bigr)\bigr)x^{s_iy}.
\end{gather*}
Substituting into~\eqref{spointred} and using the formulas~\eqref{Aisum} and~\eqref{Aires}, we now also obtain~\eqref{formula} in case $\alpha_i\smash{\bigl(Q^\vee\bigr)}=\mathbb{Z}$.


The next step is choosing the normalization factors $\kappa_w^\mathcal{O}$, $w\in W_0$, in such a way that the coefficient of $x^{s_i\lambda}$ in~\eqref{formula}
only depends on $y=\lambda+wc^\mathcal{O}$. We will show that this is the case for the scalars
\begin{equation}\label{kappawy}
\kappa_w^\mathcal{O}:=\prod_{\alpha\in\Pi(w)}\mathbf{k}_\alpha^{\chi_e(\alpha(c^\mathcal{O}))}\mathbf{k}_{\alpha/2}^{-\chi_o(\alpha(c^\mathcal{O}))}
\end{equation}
(a small computation using~\eqref{W0C} shows that $\kappa_w^\mathcal{O}$ reduces for $\mathbf{k}\in\mathcal{K}^{{\rm res}}$ to the normalization factor~\cite[(4.11)]{SSV2}).
Since $\Pi(w)=\Phi_0^+\cap w^{-1}\Phi_0^-$ for $w\in W_0$, we have
\begin{equation*}
\Pi(s_iw)=
\begin{cases}
\Pi(w)\cup\big\{w^{-1}\alpha_i\big\} &\text{if}\ \chi_+\bigl(w^{-1}\alpha_i\bigr)=1,\\
\Pi(w)\setminus\big\{-w^{-1}\alpha_i\big\} &\text{if} \
\chi_-\bigl(w^{-1}\alpha_i\bigr)=1.
\end{cases}
\end{equation*}
Combined with the fact that $\chi_e,\chi_o\colon \mathbb{R}\rightarrow\{0,1\}$ are even functions, we obtain
\begin{equation*}
\frac{\kappa_{s_iw}^\mathcal{O}}{\kappa_w^\mathcal{O}}=
\begin{cases}
k_i^{\chi_e(\alpha_i (wc^\mathcal{O}))}u_i^{-\chi_o(\alpha_i (wc^\mathcal{O} ))} &\text{if}\ \chi_+\bigl(w^{-1}\alpha_i\bigr)=1,\\
k_i^{-\chi_e(\alpha_i (wc^\mathcal{O}))}u_i^{\chi_o(\alpha_i (wc^\mathcal{O} ))} &\text{if}\ \chi_-\bigl(w^{-1}\alpha_i\bigr)=1,
\end{cases}
\end{equation*}
and hence
\begin{gather*}%\label{claim}
\frac{\kappa_{s_iw}^\mathcal{O}}{\kappa_w^\mathcal{O}}
+\bigl(k_i-k_i^{-1}\bigr)\chi_-\bigl(w^{-1}\alpha_i\bigr)\chi_e\bigl(\alpha_i\bigl(wc^\mathcal{O}\bigr)\bigr)+\bigl(u_i-u_i^{-1}\bigr)\chi_+\bigl(w^{-1}\alpha_i\bigr)\chi_o\bigl(\alpha_i\bigl(wc^\mathcal{O}\bigr)\bigr)\\
\qquad =k_i^{\chi_e(\alpha_i\bigl(wc^\mathcal{O}\bigr))}u_i^{\chi_o(\alpha_i\bigl(wc^\mathcal{O}\bigr))}
\end{gather*}
for $w\in W_0$ and $i=1,\dots,r$. So formula~\eqref{formula} reduces
for the specific choice~\eqref{kappawy} of $\kappa_w^\mathcal{O}$ to
\begin{align}
\frac{1}{\kappa_w^\mathcal{O}}\psi^\mathcal{O}\bigl(T_ix^\lambda T_w\bigr)
={}&k_i^{\chi_e(\alpha_i (wc^\mathcal{O} ))}u_i^{\chi_o(\alpha_i (wc^\mathcal{O} ))}x^{s_iy}\nonumber\\
&{}{+}\ \bigl(k_i-k_i^{-1}\bigr)\nabla_i^e(x^y)+\bigl(u_i-u_i^{-1}\bigr)\nabla_i^o(x^y).\label{formula2}
\end{align}
Now note that the map
\begin{equation}\label{functionQ}
E\rightarrow\mathbf{F}^\times,\qquad y\mapsto k_i^{\chi_e(\alpha_i(y))}u_i^{\chi_o(\alpha_i(y))}
\end{equation}
is $\tau\smash{\bigl(Q^\vee\bigr)}$-invariant. This is trivial when $\Phi_0$ is of type ${\rm C}_r$, $r\geq 1$, and $i=r$, since in this case $\alpha_r\smash{\bigl(Q^\vee\bigr)}=\mathbb{Z}_e$. In all other cases, $\alpha_i\smash{\bigl(Q^\vee\bigr)}=\mathbb{Z}$ hence $k_i=u_i$, in which case it follows from the fact that the function~\eqref{functionQ} reduces to $y\mapsto k_i^{\chi_{\mathbb{Z}}(\alpha_i(y))}$.
In~\eqref{formula2}, we may thus replace the coefficient
\[
k_i^{\chi_e(\alpha_i (wc^\mathcal{O} ))}u_i^{\chi_o(\alpha_i (wc^\mathcal{O} ))}
\] of $x^{s_iy}$ by
\smash{$k_i^{\chi_e(\alpha_i(y))}u_i^{\chi_o(\alpha_i(y))}$}.

In conclusion, for $\kappa_w^\mathcal{O}$ given by~\eqref{kappawy} the associated linear map $\psi^\mathcal{O}$~\eqref{psicmap} satisfies
\[
\psi^\mathcal{O}(T_ih)=\mathcal{D}_i^\mathcal{O}\bigl(\psi^\mathcal{O}(h)\bigr)\qquad \forall h\in H^X
\]
with $\mathcal{D}_i^\mathcal{O}\in\textup{End}(\mathbf{F}[\mathcal{O}])$ defined by
\begin{equation*}
\mathcal{D}_i^\mathcal{O}(x^y):=k_i^{\chi_e(\alpha_i(y))}u_i^{\chi_o(\alpha_i(y))}x^{s_iy}+\bigl(k_i-k_i^{-1}\bigr)\nabla_i^e(x^y)+\bigl(u_i-u_i^{-1}\bigr)\nabla_i^o(x^y)
\end{equation*}
for $y\in\mathcal{O}$.
The kernel of $\psi^\mathcal{O}\colon H^X\twoheadrightarrow\mathbf{F}[\mathcal{O}]$ thus is a left-ideal, and $\mathbf{F}[\mathcal{O}]$ inherits a $H^X$-action from $H^X/\textup{ker}\bigl(\psi^\mathcal{O}\bigr)$ with $T_i$ acting by $\mathcal{D}_i^\mathcal{O}=\pi^\mathcal{O}(T_i)$ and $x^\mu$ acting by
$\pi^\mathcal{O}(x^\mu)$. This completes the proof of the theorem.
\end{proof}

\begin{Remark}\label{H_remark}
The proof of Theorem~\ref{theoremH} involves a particular choice of normalisation factors $\kappa_w^{\mathcal{O}}\in\mathbf{F}^\times$ ($w\in W_0$). Any choice of $\kappa_w^\mathcal{O}$, $w\in W_0$, such that the coefficient of $x^{s_iy}$, ${y=\lambda+wc^{\mathcal{O}}}$, in~\eqref{formula} only depends on the coset $w\{v\in W_0 \mid vc^\mathcal{O}=c^\mathcal{O}\}$ for all ${y=\lambda+wc^\mathcal{O}\in\mathcal{O}}$ and ${i\in\{1,\dots,r\}}$
will lead to an explicit $H^X$-representation on $\mathbf{F}[\mathcal{O}]$ involving truncated Demazure--Lusztig type operators. The present choice~\eqref{kappawy} corresponds to a natural class of parabolically induced $H^X$-modules,
see Section~\ref{section_pi_mod} for details.
\end{Remark}
%%%%%%%%%%%%%%%%%%%%%%%%%%%%%%%%%%%%%%%%%%%
\begin{Corollary}\label{corpsi}\hfill
\begin{enumerate}\itemsep=0pt
\item[$(1)$] $\mathbf{F}[E]=\bigoplus_{\mathcal{O}\in E/W}\mathbf{F}[\mathcal{O}]$ is a decomposition of $\mathbf{F}[E]$ in $H^X$-submodules.
\item[$(2)$] Let $\mathcal{O}\in E/W$. Then
\begin{equation*}%\label{pimonomialformula}
\pi\bigl(x^\lambda T_w\bigr)x^{c^\mathcal{O}}=\kappa_w^\mathcal{O}x^{\lambda+wc^\mathcal{O}}
\end{equation*}
for $\lambda\in Q^\vee$ and $w\in W_0$,
with $\kappa_w^\mathcal{O}$ defined by~\eqref{kappawy}.
\end{enumerate}
\end{Corollary}

 \begin{proof}
 (1) This was remarked in the first paragraph of the proof of Theorem~\ref{theoremH}.

(2) By the last paragraph of the proof of Theorem~\ref{theoremH}, the epimorphism $\psi^\mathcal{O}\colon H^X\twoheadrightarrow\mathbf{F}[\mathcal{O}]$, mapping
$x^\lambda T_w$ to $\kappa_w^\mathcal{O}x^{\lambda+wc^\mathcal{O}}$ for $\lambda\in Q^\vee$ and $w\in W_0$, is $H^X$-linear for the special choice~\eqref{kappawy} of the scalars $\kappa_w^\mathcal{O}$. Hence
\[
\pi\bigl(x^\lambda T_w\bigr)x^{c^\mathcal{O}}=\pi\bigl(x^\lambda T_w\bigr)\psi^\mathcal{O}(1)=\psi^\mathcal{O}\bigl(x^\lambda T_w\bigr)=
\kappa_w^\mathcal{O}x^{\lambda+wc^\mathcal{O}}
\]
for $\lambda\in Q^\vee$ and $w\in W_0$.
 \end{proof}

 For the upgrade of Theorem~\ref{theoremH} to the double affine Hecke algebra $\mathbb{H}$ (see Section~\ref{wtsection}), it is useful to introduce the notation
 \begin{equation*}%\label{piO2}
 \pi^\mathcal{O}(\cdot):=\pi(\cdot)\vert_{\mathbf{F}[\mathcal{O}]}
 \end{equation*}
 for the representation map of the $H^X$-submodule $\mathbf{F}[\mathcal{O}]$ (as we in fact already have done in the proof of Theorem~\ref{theoremH}). The reason for this is that the extension of $\pi^{\mathcal{O}}\colon H^X\rightarrow\textup{End}(\mathbf{F}[\mathcal{O}])$ to a~$\mathbb{H}$-representation on $\mathbf{F}[\mathcal{O}]$ involves additional $\mathcal{O}$-dependent representation parameters.

\begin{Remark}\quad
\begin{enumerate}\itemsep=0pt
\item
By~\eqref{Aisum}, we have
\begin{equation*}%\label{kisucase}
\pi(T_i)x^y=k_i^{\chi_{\mathbb{Z}}(\alpha_i(y))}x^{s_iy}+\bigl(k_i-k_i^{-1}\bigr)\nabla_i(x^y)\qquad
\text{if}\ k_i=u_i,
\end{equation*}
from which it follows that the $H^X$-representation $\pi^\mathcal{O}$ for $\mathbf{k}\in\mathcal{K}^{{\rm res}}$ is the restriction to $H^X$ of the quasi-polynomial $\mathbb{H}$-representation defined in~\cite[Theorem~1.1] {SSV2}.
\item Let $\Lambda\subset E$ be a $W_0$-invariant lattice containing $Q^\vee$. Then $\mathbf{F}[\Lambda]$ is a $\pi\bigl(H^X\bigr)$-submodule, and the action of $\pi(T_i)\vert_{\mathbf{F}[\Lambda]}$ can be written in terms of metaplectic Demazure--Lusztig type operators when $k_i=u_i$, see~\cite{SSV2}.
The resulting \smash{$H^X$}-representation $\pi(\cdot)\vert_{\mathbf{F}[\Lambda]}$ for $\mathbf{k}\in\mathcal{K}^{{\rm res}}$ is essentially the one introduced in~\cite[Theorem~3.7]{SSV1}. The proof of
 Theorem~\ref{theoremH} basically follows the same strategy as the proof of~\cite[Theorem~3.7]{SSV1}.
 \end{enumerate}
 \end{Remark}
 %%%%%%%%%%%%%%%%%%%%%%%%%%%%%%%%%%%%%%%%%%%%
\subsection{Parabolic data and parabolically induced modules}\label{section_pi_mod}
%%%%%%%%%%%%%%%%%%%%%%%%%%%%%%%%%%%%%%%%%%%%%
We show in this subsection that the representation $\pi^\mathcal{O}$ is a parabolically induced $H^X$-module when $\alpha_0\bigl(c^\mathcal{O}\bigr)\not=0$.

For a subset $I\subseteq\{1,\dots,r\}$, write
\begin{itemize}\itemsep=0pt
\item $W_{0,I}$ for the parabolic subgroup of $W_0$ generated by $s_i$, $i\in I$,
\item $W_0^I$ for the minimal coset representatives of $W_0/W_{0,I}$,
\item $H_{0,I}$ for the subalgebra of $H_0$ generated by $T_i$, $i\in I$.
\end{itemize}
The finite Hecke algebra $H_0$ is a free right $H_{0,I}$-module with basis $\{T_v\}_{v\in W_0^I}$, since
\begin{equation}\label{lengthadd}
\ell(vw)=\ell(v)+\ell(w)\qquad \text{for}\ v\in W_0^I \ \text{and}\ w\in W_{0,I}.
\end{equation}
Here is an immediate lift to the affine Hecke algebra $H^X$.
%%%%%%%%%%%%%%%%%%%%%%%%%%%%%%%%%%%%%%%%%%%%%
\begin{Lemma}\label{lemfree}\hfill
\begin{enumerate}\itemsep=0pt
\item[$(1)$]
$\{\tau(\lambda)v\}_{(\lambda,v)\in Q^\vee\times W_0^I}$ is a complete set of representatives of $W/W_{0,I}$.
\item[$(2)$] $H^X$ is a free right $H_{0,I}$-module with
basis \smash{$\bigl\{x^\lambda T_v\bigr\}_{(\lambda,v)\in Q^\vee\times W_0^I}$}.
\end{enumerate}
\end{Lemma}
%%%%%%%%%%%%%%%%%%%%%%%%%%%%%%%%%%%%%%%%%%%%%
\begin{proof}
The first part is a consequence of the fact that $(\lambda,w)\mapsto \tau(\lambda)w$ defines a bijection $Q^\vee\rtimes W_0\overset{\sim}{\longrightarrow} W$. For the second part, note that the multiplication map of $H^X$ restricts to
a~linear isomorphism $\mathbf{F}\smash{\bigl[Q^\vee\bigr]}\otimes H_0\overset{\sim}{\longrightarrow} H^X$.
\end{proof}

The more familiar parabolic structures on $W$ and on the associated affine Hecke algebra $H$ arise from their Coxeter type presentations. In this case it depends on a subset $J$ of $\{0,\dots,r\}$. We~write
\begin{itemize}\itemsep=0pt
\item $W_{J}$ for the parabolic subgroup of $W$ generated by $s_j$, $j\in J$,
\item $W^J$ for the minimal coset representatives of $W/W_{J}$,
\item $H_{J}$ for the subalgebra of $H$ generated by $T_j$, $j\in J$.
\end{itemize}
The length identity~\eqref{lengthadd} now also holds true for $v\in W^J$ and $w\in W_J$. As a consequence, $H$~is a free right $H_J$-module with basis
$\{T_g\}_{g\in W^J}$.

The closure $\overline{C_+}$ of the fundamental alcove $C_+$ splits in a disjoint union of facets
\begin{equation*}%\label{facetC}
\overline{C_+}=\bigsqcup_{J\subsetneq \{0,\dots,r\}}C_+^J,
\end{equation*}
with $C_+^J$ the set of vectors $y\in\overline{C_+}$ for which $\alpha_j(y)=0$ if and only if $j\in J$.
For $y\in E$ denote by $W_y\subset W$ the subgroup of $W$ fixing $y$. It is well known that
\[
W_c=W_J\qquad \text{for}\ c\in C_+^J.
\]

\begin{Definition}\label{IJ}
For a $W$-orbit $\mathcal{O}$ in $E$, we write
$J(\mathcal{O})$ for the subset of $\{0,\dots,r\}$ such that $c^\mathcal{O}\in C_+^{J(\mathcal{O})}$. We furthermore write
\begin{equation*}%\label{I(O)}
I(\mathcal{O}):=J(\mathcal{O})\cap\{1,\dots,r\}.
\end{equation*}
\end{Definition}

Note that $I(\mathcal{O})=J(\mathcal{O})$ if and only if $\alpha_0\bigl(c^\mathcal{O}\bigr)\not=0$.
We will use the shorthand notations
\[
C_+^\mathcal{O},\ nW_\mathcal{O},\ W^\mathcal{O},\ H_\mathcal{O}, \ \dots
\]
for $C_+^{J(\mathcal{O})},W_{J(\mathcal{O})}, W^{J(\mathcal{O})}, H_{J(\mathcal{O})},\dots$ and
\[
W_{0,\mathcal{O}},\ W^{\mathcal{O}}_0,\ H_{0,\mathcal{O}},\ \dots
\]
for $W_{0,I(\mathcal{O})}, W_0^{I(\mathcal{O})}, H_{0,I(\mathcal{O})},\dots$.
The following lemma refines the decomposition~\eqref{FOpoldec} of the space of quasi-polynomials $\mathbf{F}[\mathcal{O}]$ as $\mathbf{F}\smash{\bigl[Q^\vee\bigr]}$-module when
$\alpha_0\bigl(c^\mathcal{O}\bigr)\not=0$.
%%%%%%%%%%%%%%%%%%%%%%%%%%%%%
\begin{Lemma}\label{Pcpoldec}
If $\mathcal{O}$ is a $W$-orbit such that $\alpha_0\bigl(c^\mathcal{O}\bigr)\not=0$, then
$\mathbf{F}[\mathcal{O}]$ decomposes as
\begin{equation}\label{FpoldecI}
\mathbf{F}[\mathcal{O}]=\bigoplus_{v\in W_0^\mathcal{O}}\mathbf{F}\smash{\bigl[Q^\vee\bigr]}x^{vc^\mathcal{O}}.
\end{equation}
\end{Lemma}
%%%%%%%%%%%%%%%%%%%%%%%%%%%%%%
\begin{proof}
The assignment $gW_{0,\mathcal{O}}\mapsto gc^\mathcal{O}$ gives rise to a bijection
\[
W/W_{0,\mathcal{O}}\overset{\sim}{\longrightarrow}\mathcal{O},
\]
since $W_{c^\mathcal{O}}=W_{\mathcal{O}}=W_{0,\mathcal{O}}$ by the assumption on $\mathcal{O}$, and $W_0^{\mathcal{O}}$ is a complete set of representatives of the double coset space $\tau\smash{\bigl(Q^\vee\bigr)}\backslash W/W_{0,\mathcal{O}}$ by Lemma
\ref{lemfree}(1).
Hence $\{vc^\mathcal{O}\}_{v\in W_0^{\mathcal{O}}}=W_0c^\mathcal{O}$ is a~complete set of representatives of the $\tau\smash{\bigl(Q^\vee\bigr)}$-orbits in $\mathcal{O}$, and the lemma follows.
\end{proof}

For a $W$-orbit $\mathcal{O}$ in $E$, let $\mathbf{F}1^\mathcal{O}$
the trivial $H_\mathcal{O}$-module, defined by
\[
T_j1^\mathcal{O}=k_j1^\mathcal{O}\qquad \text{for}\ j\in J(\mathcal{O}).
\]
We will also view $\mathbf{F}1^\mathcal{O}$ as $H_{0,\mathcal{O}}$-module by restricting the action to $H_{0,\mathcal{O}}$.
Consider the induced $H^X$-module
\[
{\rm Ind}_{H_{0,\mathcal{O}}}^{H^X}\bigl(\mathbf{F}1^{\mathcal{O}}\bigr)=H^X\otimes_{H_{0,\mathcal{O}}}\mathbf{F}1^\mathcal{O}
\]
and write
\[
\mathbf{1}^{\mathcal{O}}=1\otimes_{H_{0,\mathcal{O}}}1^\mathcal{O}
\]
for its canonical cyclic vector.

%%%%%%%%%%%%%%
\begin{Proposition}\label{indcor}
There exists a unique $H^X$-linear epimorphism
\begin{equation}\label{epiuse}
{\rm Ind}_{H_{0,\mathcal{O}}}^{H^X}\bigl(\mathbf{F}1^{\mathcal{O}}\bigr)\twoheadrightarrow
\bigl(\mathbf{F}[\mathcal{O}],\pi^\mathcal{O}\bigr)
\end{equation}
mapping $\mathbf{1}^{\mathcal{O}}$ to \smash{$x^{c^\mathcal{O}}$}. It is an isomorphism when $\alpha_0\bigl(c^\mathcal{O}\bigr)\not=0$.
\end{Proposition}

\begin{proof}
For the first statement, we need to show that the assignment $h\mathbf{1}^{\mathcal{O}}\mapsto\pi(h)x^{c^\mathcal{O}}$ for $h\in H^X$ is well defined. It suffices to note that
\[
\pi(T_i)x^{c^\mathcal{O}}=k_ix^{c^\mathcal{O}}\qquad \text{for}\ i\in I(\mathcal{O}).
\]
But for $i\in I(\mathcal{O})$, we have \smash{$s_ic^\mathcal{O}=c^{\mathcal{O}}$}, and hence
\[
\pi(T_i)x^{c^\mathcal{O}}=\kappa_{s_i}^\mathcal{O} x^{s_ic^\mathcal{O}}=k_ix^{c^\mathcal{O}}
\]
by Corollary~\ref{corpsi}\,(2) and~\eqref{kappawy}.

For the second statement, note first that
\begin{equation}\label{basis}
\bigl\{x^\lambda T_v\mathbf{1}^{\mathcal{O}}\mid (\lambda,v)\in Q^\vee\times W_0^\mathcal{O}\bigr\}
\end{equation}
is a basis of \smash{${\rm Ind}_{H_{0,\mathcal{O}}}^{H^X}\bigl(\mathbf{F}1^\mathcal{O}\bigr)$}. By Corollary~\ref{corpsi}\,(2), the basis element $x^\lambda T_v\mathbf{1}^{\mathcal{O}}$
is mapped by the epimorphism~\eqref{epiuse} to
\[
\pi(x^\lambda T_v)x^{c^\mathcal{O}}=\kappa_v^\mathcal{O}x^{\lambda+vc^\mathcal{O}}.
\]
By Lemma~\ref{Pcpoldec}, we conclude that the epimorphism~\eqref{epiuse} maps the basis~\eqref{basis} of
${\rm Ind}_{H_{0,\mathcal{O}}}^{H^X}\bigl(\mathbf{F}1^\mathcal{O}\bigr)$
to a basis of $\mathbf{F}[\mathcal{O}]$, which concludes the proof of the second statement.
\end{proof}

The natural analog of Proposition~\ref{indcor} for $W$-orbits $\mathcal{O}$ with $\alpha_0\bigl(c^\mathcal{O}\bigr)=0$ (i.e., with $0\in J(\mathcal{O})$) requires the extension of the $H^X$-action $\pi^\mathcal{O}$ on $\mathbf{F}[\mathcal{O}]$ to an action of the double affine Hecke algebra $\mathbb{H}$. This will be the subject of the next subsection.

\subsection[The quasi-polynomial representation of H]{The quasi-polynomial representation of $\boldsymbol{\mathbb{H}}$}\label{wtsection}
%%%%%%%%%%%%%%%%%%%%%
We now promote the quasi-polynomial representation $\pi^\mathcal{O}$ of the affine Hecke algebra $H^X$ to a~family of representations of the double affine Hecke algebra $\mathbb{H}$.
The number of additional parameters
depends on the facet \smash{$C^{J(\mathcal{O)}}_+$} containing $c^\mathcal{O}$.
For $\mathbf{k}\in\mathcal{K}^{{\rm res}}$ the extended representations are the quasi-polynomial $\mathbb{H}$-representations $\pi_{c^\mathcal{O},t}$ introduced in~\cite[Theorem~1.1]{SSV2}. For~${\mathbf{k}\in\mathcal{K}}$ and~$\Phi_0$ of type ${\rm C}_r$, $r\geq 1$, they will give quasi-polynomial extensions of the polynomial representation of the
type ${\rm C}^\vee{\rm C}_r$ double affine Hecke algebra $\mathbb{H}$~\cite{No,Sa}.

Recall the definition of the $\mathbf{F}$-torus $\mathbf{T}$ (see Definition~\ref{T_def}). For $J\subsetneq\{0,\dots,r\}$, consider its affine subtori
\begin{equation*}
\mathbf{T}_J:=\bigl\{t\in\mathbf{T}\mid t^{\alpha_j^\vee}=1\ \forall j\in J\bigr\}.
\end{equation*}
For a $W$-orbit $\mathcal{O}$ in $E$, write
\[\mathbf{T}_\mathcal{O}:=\mathbf{T}_{J(\mathcal{O})}.
\]
Note that
\begin{equation}\label{sat}
s_at=t\bigl(t^{-a^\vee}\bigr)^{\overline{a}}\qquad \text{for}\ a\in\Phi\ \text{and}\ t\in\mathbf{T},
\end{equation}
where \smash{$t\bigl(t^{-a^\vee}\bigr)^{\overline{a}}\in\mathbf{T}$} is viewed as character of $Q^\vee$ by \smash{$\lambda\mapsto t^\lambda\bigl(t^{-a^\vee}\bigr)^{\overline{a}(\lambda)}$} for
$\lambda\in Q^\vee$. By~\eqref{sat}, we~have
$s_at=t$ if $t^{a^\vee}=1$, and so
\begin{equation}\label{containedinv}
\mathbf{T}_J\subseteq\mathbf{T}^{W_J}:=\{t\in\mathbf{T} \mid gt=t\ \forall g\in W_J\}.
\end{equation}
Note that $\mathbf{T}_J$ and $\mathbf{T}^{W_J}$ are sub-tori of $\mathbf{T}$ when $0\not\in J$.

In~\cite[Lemma 4.2]{SSV2}, the $W_0$-action on $\mathbf{F}[\mathcal{O}]$ was extended to a family of $W$-actions on $\mathbf{F}[\mathcal{O}]$, parametrised by
$t\in\mathbf{T}_\mathcal{O}$. These actions are compatible with the $W$-action~\eqref{Wa} on $\mathbf{F}\smash{\bigl[Q^\vee\bigr]}$ by $q$-dilations and reflections.
The definition of this action requires the following definition.

\begin{Definition}\label{gy}
For $y\in E$, write $\mathbf{g}_y\in W$ for the unique element of shortest length in $W$ such that $\mathbf{g}_y^{-1}y\in\overline{C_+}$.
\end{Definition}
%%%%%%%%%%%%%%
Note that if $y$ lies in the $W$-orbit $\mathcal{O}$ of $E$,
then $\mathbf{g}_y$ is the unique element in $W^\mathcal{O}$ such that
$y=\mathbf{g}_yc^\mathcal{O}$.
%%%%%%%%%%%%%%%%%%%%%%%%%%
\begin{Lemma}\label{tdefaction}
Let $\mathcal{O}$ be a $W$-orbit in $E$. For $t\in\mathbf{T}_\mathcal{O}$,
 the formulas
\begin{gather}
w_tx^y:=w(x^y)=x^{wy},\qquad w\in W_0, \nonumber\\
\tau(\lambda)_tx^y:=(\mathbf{g}_yt)^{-\lambda}x^y,\qquad \lambda\in Q^\vee\label{taction}
\end{gather}
for $y\in\mathcal{O}$ define a linear left $W$-action on $\mathbf{F}[\mathcal{O}]$
 satisfying
\begin{equation}\label{wcomp}
g_t(pf)=(gp)(g_tf)
\end{equation}
for $g\in W$, $p\in\mathbf{F}\smash{\bigl[Q^\vee\bigr]}$ and $f\in\mathbf{F}[\mathcal{O}]$.
\end{Lemma}
%%%%%%%%%%%%%%%%%%%%%%%%%%%
It is instructive to recall the proof of~\eqref{wcomp} for $g=\tau(\lambda)$, $\lambda\in Q^\vee$. First note that by~\eqref{containedinv},
we may replace $\mathbf{g}_y$ in~\eqref{taction} by any other representative of the coset $\mathbf{g}_yW_\mathcal{O}$.
We then have for~${\mu\in Q^\vee}$,
\begin{align}
\tau(\lambda)_tx^{\mu+y}&=(\mathbf{g}_{\mu+y}t)^{-\lambda}x^{\mu+y} \nonumber\\
&=(\tau(\mu)\mathbf{g}_yt)^{-\lambda}x^{\mu+y} \nonumber\\
&=q^{-\langle\lambda,\mu\rangle}(\mathbf{g}_y t)^{-\lambda}x^{\mu+y} \nonumber\\
&=\bigl(q^{-\langle\lambda,\mu\rangle}x^\mu\bigr)\bigl((\mathbf{g}_yt)^{-\lambda}x^y\bigr),\label{tcomp}
\end{align}
where we used $\mathbf{g}_{y+\mu}W_\mathcal{O}=\tau(\mu)\mathbf{g}_yW_\mathcal{O}$ for the second equality, and~\eqref{Waction} for the third equality.
The last line in~\eqref{tcomp} clearly equals $(\tau(\lambda)x^\mu)(\tau(\lambda)_tx^y)$. Note in particular that the family of $W$-actions~\eqref{taction} depends on $q$ via the $W$-action~\eqref{Waction} on $\mathbf{T}$.

Note that by~\eqref{tcomp} we have the formula
\[
\tau(\lambda)_t\vert_{\mathbf{F}[Q^\vee]x^y}=(\mathbf{g}_yt)^{-\lambda}(x^y\circ\tau(\lambda)\circ x^{-y})\vert_{\mathbf{F}[Q^\vee]x^y},
\]
where $x^{\pm y}$ are regarded as multiplication operators on $\mathbf{F}[E]$.

The following result extends~\cite[Theorem~1.1]{SSV2} to multiplicity functions $\mathbf{k}\in\mathcal{K}$.
%%%%%%%%%%%%%%%%%%%%%%%%%%%%%%%%%%%%%%%%%%%%%%%%%%%%%
\begin{Theorem}\label{mthm}
Let $\mathcal{O}$ be a $W$-orbit in $E$ and $t\in\mathbf{T}_\mathcal{O}$. The formulas
\begin{gather}
\pi_{t}^\mathcal{O}\bigl(x^\lambda\bigr)x^y:=x^{y+\lambda}, \nonumber\\
\pi_{t}^\mathcal{O}(T_j)x^y:=k_j^{\chi_e(\overline{\alpha}_j(y))}u_j^{\chi_o(\overline{\alpha}_j(y))}s_{j,t}x^y+\bigl(k_j-k_j^{-1}\bigr)\nabla_j^e(x^y)+
\bigl(u_j-u_j^{-1}\bigr)\nabla_j^o(x^y)
\label{qprepHH}
\end{gather}
for $j=0,\dots,r$, $\lambda\in Q^\vee$ and $y\in\mathcal{O}$ define a representation $\pi_{t}^\mathcal{O}\colon \mathbb{H}\rightarrow\textup{End}(\mathbf{F}[\mathcal{O}])$.
\end{Theorem}
%%%%%%%%%%%%%%%
\begin{Remark}
Note that $\pi_{t}^\mathcal{O}\vert_{H^X}$ is the restriction $\pi^{\mathcal{O}}$ of the quasi-polynomial $H^X$-represen\-tation $\pi$ from Theorem~\ref{theoremH} to $\mathbf{F}[\mathcal{O}]$. Furthermore, for restricted multiplicity
parameters $\mathbf{k}\in\mathcal{K}^{{\rm res}}$ we have
\[
\pi_{t}^\mathcal{O}(T_j)=k_j^{\chi_{\mathbb{Z}}(\overline{\alpha}_j(y))}s_{j,t}x^y+\bigl(k_j-k_j^{-1}\bigr)\nabla_j(x^y)
\]
since $k_j=u_j$, hence $\pi_{t}^\mathcal{O}$ then coincides with the quasi-polynomial representation $\pi_{c^\mathcal{O},t}$ defined in~\cite[Theorem~1.1]{SSV2}.
\end{Remark}
%%%%%%%%%%%%%%%
\begin{proof}
Consider $\mathbb{H}$ as left regular $\mathbb{H}$-module.
Relative
to the $\mathbf{F}$-basis
\[
\bigl\{x^\lambda T_w Y^\mu\mid \lambda,\mu\in Q^\vee,\, w\in W_0\bigr\}
\]
of $\mathbb{H}$, the $H$-action is explicitly given by
\begin{gather}
T_jx^\lambda T_wY^\mu=s_j\bigl(x^\lambda\bigr)T_{s_{\overline{\alpha}_j}w}Y^{\mu-w^{-1}s_j(0)} \nonumber\\
\hphantom{T_jx^\lambda T_wY^\mu=}{}
+\big(k_j-k_j^{-1}\big)\Biggl(\frac{1-x^{(2\chi_-(w^{-1}\overline{\alpha}_j)-\overline{\alpha}_j(\lambda))\alpha_j^\vee}}{1-x^{2\alpha_j^\vee}}\Biggr)x^\lambda T_wY^\mu\nonumber\\
\hphantom{T_jx^\lambda T_wY^\mu=}{}
+\big(u_j-u_j^{-1}\big)\Biggl(\frac{x^{\alpha_j^\vee}-x^{(1-\overline{\alpha}_j(\lambda))\alpha_j^\vee}}{1-x^{2\alpha_j^\vee}}\Biggr)x^\lambda T_w Y^\mu\label{regularHH}
\end{gather}
for $w\in W_0$, $\lambda,\mu\in Q^\vee$ and $j=0,\dots,r$. For $j\in\{1,\dots,r\}$, formula~\eqref{regularHH} follows immediately from~\eqref{regularHX}.
For $j=0$, formula~\eqref{regularHH} follows by commuting $T_0$ and $x^\lambda$ using the cross relation~\eqref{crossX} and then applying
the identity
\begin{equation}\label{T0Tw}
T_0T_w=T_{s_{\overline{\alpha}_0}w}Y^{-w^{-1}s_0(0)}+\chi_-\bigl(w^{-1}\overline{\alpha}_0\bigr)\bigl(k_0-k_0^{-1}\bigr)T_w
\end{equation}
in $H$. For the proof of~\eqref{T0Tw}, first note that it is equivalent to the identity
\begin{equation}\label{Hid}
T_0^{\chi(w^{-1}\varphi)}T_{s_\varphi w}=T_wY^{w^{-1}\varphi^\vee}
\end{equation}
in $H$ since $s_0(0)=\varphi^\vee$, $\overline{\alpha}_0=-\varphi$ and $T_0^{-1}=T_0-k_0+k_0^{-1}$.
For a proof of~\eqref{Hid} see, for instance,~\cite[(3.3.6)]{Ma}.

Let $\mathcal{O}$ be a $W$-orbit in $E$.
Recall the linear epimorphism $\psi^\mathcal{O}\colon H^X\twoheadrightarrow\mathbf{F}[\mathcal{O}]$,
defined by~\eqref{psicmap}, which is an epimorphism of $H^X$-modules by Corollary~\ref{corpsi}\,(2). The goal is to find extensions of~$\psi^{\mathcal{O}}$ to linear epimorphisms $\mathbb{H}\twoheadrightarrow\mathbf{F}[\mathcal{O}]$ such that
\begin{enumerate}\itemsep=0pt
\item[(1)] their kernels are left ideals in $\mathbb{H}$,
\item[(2)]
right multiplication by $\mathbf{F}_Y\smash{\bigl[Q^\vee\bigr]}$ is turned into multiplication by a linear character of $\mathbf{F}_Y\smash{\bigl[Q^\vee\bigr]}$.
\end{enumerate}
This will upgrade the $H^X$-action on $\mathbf{F}[\mathcal{O}]$ to a family of $\mathbb{H}$-actions satisfying the additional property that the quasi-monomial \smash{$x^{c^\mathcal{O}}$} will be a simultaneous eigenvector
for the action of $\mathbf{F}_Y\smash{\bigl[Q^\vee\bigr]}$. The family of extended $\mathbb{H}$-actions on $\mathbf{F}[\mathcal{O}]$ will be natural parametrised by the associated linear characters of $\mathbf{F}_Y\smash{\bigl[Q^\vee\bigr]}$, which in turn can be described by an affine subtorus of the form $\mathfrak{s}\mathbf{T}_{\mathcal{O}}$ for a specific basepoint $\mathfrak{s}\in\mathbf{T}$, so be determined in due course.

So our starting point will be the desired property (2). Fix $t\in\mathbf{T}_{\mathcal{O}}$ and consider the extension of $\psi^\mathcal{O}$ to a ($t$-dependent) linear map $\psi_{t}^\mathcal{O}\colon \mathbb{H}\twoheadrightarrow\mathbf{F}[\mathcal{O}]$ by
\begin{equation}\label{psict}
\psi_{t}^\mathcal{O}\bigl(x^\lambda T_wY^\mu\bigr):=(\mathfrak{s}t)^{-\mu}\psi^\mathcal{O}\bigl(x^\lambda T_w\bigr)=\kappa_w^\mathcal{O}(\mathfrak{s}t)^{-\mu}x^{\lambda+wc^\mathcal{O}}
\end{equation}
for $w\in W_0$ and $\lambda,\mu\in Q^\vee$ \big(recall here that $\kappa_w^\mathcal{O}\in\mathbf{F}^\times$ is the explicit scalar defined by~\eqref{kappawy}\big).
By the first formula of~\eqref{psict} and Corollary~\ref{corpsi}\,(2), it follows that
\[
\psi_{t}^\mathcal{O}(hh^\prime)
=\pi_{t}^\mathcal{O}(h)\psi_{t}^\mathcal{O}(h^\prime)\qquad \text{for}\, h\in H^X\, \text{and}\, h^\prime\in\mathbb{H}.
\]
In particular, the kernel of $\psi_{t}^\mathcal{O}$ is a left $H^X$-submodule in $\mathbb{H}$.
To meet the first property (1), we~now fine-tune the choice of $\mathfrak{s}\in\mathbf{T}$ such that
\[
\psi_{t}^{\mathcal{O}}(T_0h^\prime)=\mathcal{D}_0\bigl(\psi_{t}^\mathcal{O}(h^\prime)\bigr) \qquad \textup{ for all }\, h^\prime\in\mathbb{H}
\]
for some $\mathcal{D}_0\in\textup{End}(\mathbf{F}[\mathcal{O}])$.

Note that for $\lambda\in Q^\vee$ and $w\in W_0$,
\begin{equation*}
s_{0,t}x^{\lambda+wc^\mathcal{O}}=s_0\bigl(x^\lambda\bigr)\bigl(s_{0,t}x^{wc^\mathcal{O}}\bigr)=t^{w^{-1}\varphi^\vee}s_0\bigl(x^\lambda\bigr)x^{s_\varphi wc^\mathcal{O}}
\end{equation*}
and hence, by~\eqref{regularHH},
\begin{gather}
\frac{(\mathfrak{s}t)^\mu}{\kappa_w^\mathcal{O}}\psi_{t}^\mathcal{O}(T_0x^\lambda T_wY^\mu)=\frac{\kappa_{s_\varphi w}^\mathcal{O}\,\mathfrak{s}^{w^{-1}\varphi^\vee}}{\kappa_w^\mathcal{O}}s_{0,t}x^{\lambda+wc^\mathcal{O}}\nonumber\\
\hphantom{\frac{(\mathfrak{s}t)^\mu}{\kappa_w^\mathcal{O}}\psi_{t}^\mathcal{O}(T_0x^\lambda T_wY^\mu)=}{}
+(k_0-k_0^{-1})\Biggl(\frac{1-x^{(2\chi_-(w^{-1}\overline{\alpha}_0)-\overline{\alpha}_0(\lambda))\alpha_0^\vee}}{1-x^{2\alpha_0^\vee}}\Biggr)x^{\lambda+wc^\mathcal{O}}\nonumber\\
\hphantom{\frac{(\mathfrak{s}t)^\mu}{\kappa_w^\mathcal{O}}\psi_{t}^\mathcal{O}(T_0x^\lambda T_wY^\mu)=}{}
+(u_0-u_0^{-1})\Biggl(\frac{x^{\alpha_0^\vee}-x^{(1-\overline{\alpha}_0(\lambda))\alpha_0^\vee}}{1-x^{2\alpha_0^\vee}}\Biggr)x^{\lambda+wc^\mathcal{O}}.
\label{spoint0}
\end{gather}
So we need to fine-tune $\mathfrak{s}\in\mathbf{T}$ such that the right-hand side of~\eqref{spoint0} can be written
$\mathcal{D}_0\bigl(x^{\lambda+wc}\bigr)$ for
some $\mathcal{D}_0\in\textup{End}(\mathbf{F}[\mathcal{O}])$.
We follow the proof of Theorem~\ref{theoremH}, but it requires some necessary additional computations due to the presence of the additional parameters $\mathfrak{s}$ and $t$ (compare also with~\cite[Section~5]{SSV2}, which deals with the case that $\mathbf{k}\in\mathcal{K}^{{\rm res}}$).

We first consider the case that $\Phi_0$ is of type ${\rm C}_r$, $r\geq 1$, so that $\alpha_0\smash{\bigl(Q^\vee\bigr)}=\mathbb{Z}_o$.
Fix $\lambda\in Q^\vee$ and $w\in W_0$ and set
\[
y:=\lambda+wc^\mathcal{O}.
\]
The second line of~\eqref{spoint0} can then be rewritten using the formula
\begin{equation}\label{claim3}
\Biggl(\frac{1-x^{(2\chi_-(w^{-1}\overline{\alpha}_0)-\overline{\alpha}_0(\lambda))\alpha_0^\vee}}{1-x^{2\alpha_0^\vee}}\Biggr)x^{y}=
\nabla_0^e(x^y)
+\chi_-\bigl(w^{-1}\overline{\alpha}_0\bigr)\chi_e\bigl(\overline{\alpha}_0\bigl(wc^\mathcal{O}\bigr)\bigr)s_{0,t}x^{y}.
\end{equation}
To prove~\eqref{claim3}, note that by~\eqref{W0C} we have
\begin{equation}
\label{sstep10}
2\chi_-\bigl(w^{-1}\overline{\alpha}_0\bigr)-\overline{\alpha}_0(\lambda)=
\begin{cases}
2-\lfloor\overline{\alpha}_0(y)\rfloor_e &\text{if}\ w^{-1}\overline{\alpha}_0\in\Phi_0^-\ \text{and}\ \overline{\alpha}_0\bigl(wc^\mathcal{O}\bigr)=0,\\
-\lfloor\overline{\alpha}_0(y)\rfloor_e &\text{otherwise},
\end{cases}
\end{equation}
and that $\overline{\alpha}_0\bigl(wc^\mathcal{O}\bigr)=0$ if and only if $\chi_e\bigl(\overline{\alpha}_0\bigl(wc^\mathcal{O}\bigr)\bigr)=1$.
This immediately implies~\eqref{claim3} unless
$w^{-1}\overline{\alpha}_0\in\Phi_0^-$ and $\overline{\alpha}_0\bigl(wc^\mathcal{O}\bigr)=0$.
So suppose now that $w^{-1}\overline{\alpha}_0\in\Phi_0^-$ and $\overline{\alpha}_0\bigl(wc^\mathcal{O}\bigr)=0$. Then $s_\varphi wc^\mathcal{O}=wc^\mathcal{O}$, and hence
\[
s_{0,t}x^y=s_0\bigl(x^\lambda\bigr)\bigl(s_{0,t}x^{wc^\mathcal{O}}\bigr)=\bigl(x^{\lambda-\overline{\alpha}_0(\lambda)\alpha_0^\vee}\bigr)\bigl(t^{w^{-1}\varphi^\vee}x^{wc^\mathcal{O}}\bigr)
=t^{w^{-1}\varphi^\vee}x^{\lambda-\overline{\alpha}_0(\lambda)\alpha_0^\vee+wc^\mathcal{O}}.
\]
But \smash{$\bigl(w^{-1}\varphi\bigr)\bigl(c^\mathcal{O}\bigr)=0$}, hence
\[
w^{-1}\varphi^\vee\in\Phi_0^\vee\cap\bigoplus_{i\in I(\mathcal{O})}\mathbb{Z}\alpha_i^\vee.
\]
 Since $t\in\mathbf{T}_\mathcal{O}$,
we conclude that ${t^{w^{-1}\varphi^\vee}=1}$, so
\begin{equation}\label{s0simple}
s_{0,t}x^y=x^{\lambda-\overline{\alpha}_0(\lambda)\alpha_0^\vee+wc^\mathcal{O}}.
\end{equation}
Then~\eqref{claim3} follows by combining the first case of~\eqref{sstep10} and~\eqref{s0simple}.

Similarly, we rewrite the third line of~\eqref{spoint0} using the formula
\begin{equation}\label{claim4}
\Biggl(\frac{x^{\alpha_0^\vee}-x^{(1-\overline{\alpha}_0(\lambda))\alpha_0^\vee}}{1-x^{2\alpha_0^\vee}}\Biggr)x^y=
\nabla_0^o(x^y)+
\chi_+\bigl(w^{-1}\overline{\alpha}_0\bigr)\chi_o\bigl(\overline{\alpha}_0\bigl(wc^\mathcal{O}\bigr)\bigr)
s_{0,t}x^y.
\end{equation}
For the proof of~\eqref{claim4}, we now use that
\begin{equation}\label{sstep20}
1-\overline{\alpha}_0(\lambda)=
\begin{cases}
2-\lfloor\overline{\alpha}_0(y)\rfloor_o &\text{if}\ \overline{\alpha}_0\bigl(wc^\mathcal{O}\bigr)=1,\\
-\lfloor\overline{\alpha}_0(y)\rfloor_o &\text{otherwise},
\end{cases}
\end{equation}
and the observation that $\overline{\alpha}_0\bigl(wc^\mathcal{O}\bigr)=1$ is equivalent to
$\chi_+\bigl(w^{-1}\overline{\alpha}_0\bigr)\chi_o\bigl(\overline{\alpha}_0\bigl(wc^\mathcal{O}\bigr)\bigr)\!=\!1$. Then~\eqref{claim4} is immediate if $\overline{\alpha}_0\bigl(wc^\mathcal{O}\bigr)\not=1$.
So suppose now that $\overline{\alpha}_0\bigl(wc^\mathcal{O}\bigr)=1$. Then
\begin{align}
s_{0,t}x^y
&=t^{w^{-1}\varphi^\vee}x^{\lambda-\overline{\alpha}_0(\lambda)\alpha_0^\vee+s_\varphi wc^\mathcal{O}}\nonumber\\
&=q_\varphi t^{w^{-1}\varphi^\vee}x^{\lambda-(1+\overline{\alpha}_0(\lambda))\alpha_0^\vee+wc^\mathcal{O}}=
x^{\lambda-(1+\overline{\alpha}_0(\lambda))\alpha_0^\vee+wc^\mathcal{O}}.\label{s0help}
\end{align}
The last equality follows from the fact that the affine root $a:=\bigl(w^{-1}\varphi,1\bigr)$ satisfies $a\bigl(c^\mathcal{O}\bigr)=0$, so \smash{$a^\vee\in\Phi^+\cap\bigoplus_{j\in J(\mathcal{O})}\mathbb{Z}\alpha_j^\vee$} and hence
\smash{$q_\varphi t^{w^{-1}\varphi^\vee}=t^{a^\vee}=1$} (cf.~\cite[Lemma 3.4]{SSV2}). Then~\eqref{claim4} follows from the first case of~\eqref{sstep20} and~\eqref{s0help}.


Returning now to~\eqref{spoint0},
we conclude from~\eqref{claim3} and~\eqref{claim4} that
\begin{equation}\label{formula0}
\frac{(\mathfrak{s}t)^\mu}{\kappa_w^\mathcal{O}}\psi_{t}^\mathcal{O}\bigl(T_0x^\lambda T_wY^\mu\bigr)
={\rm coeff}_w^\mathcal{O}s_{0,t}x^y+\bigl(k_0-k_0^{-1}\bigr)\nabla_0^e(x^y)+\bigl(u_0-u_0^{-1}\bigr)\nabla_0^o(x^y)
\end{equation}
with
\begin{gather}
{\rm coeff}_w^\mathcal{O}:=\frac{\kappa_{s_\varphi w}^\mathcal{O}\mathfrak{s}^{w^{-1}\varphi^\vee}}{\kappa_w^\mathcal{O}}+\bigl(k_0-k_0^{-1}\bigr)\chi_-(w^{-1}\overline{\alpha}_0)\chi_e \bigl(\overline{\alpha}_0\bigl(wc^\mathcal{O}\bigr)\bigr)\nonumber\\
\hphantom{{\rm coeff}_w^\mathcal{O}:=}{}
+\bigl(u_0-u_0^{-1}\bigr)\chi_+\bigl(w^{-1}\overline{\alpha}_0\bigr)\chi_o\bigl(\overline{\alpha}_0\bigl(wc^\mathcal{O}\bigr)\bigr).
\label{coeff}
\end{gather}
Formula~\eqref{formula0} also holds true when $\Phi_0$ is not of type ${\rm C}_r$. In fact, if $\Phi_0$ is not of type ${\rm C}_r$ then $\mathcal{K}=\mathcal{K}^{{\rm res}}$ and $k_0=u_0$,
and the formula can be recovered from~\cite[Section~5]{SSV2}. It can also be derived directly, similarly as the proof of~\eqref{formula} when $\alpha_i\smash{\bigl(Q^\vee\bigr)}=\mathbb{Z}$.
So we will now continue the proof of the theorem for $\Phi_0$ of any type, taking~\eqref{formula0} as the starting point.

Properties of the normalization factors $\kappa_w^\mathcal{O}$~\eqref{kappawy}
were derived in~\cite{SSV2} for restricted parameters $\mathbf{k}\in\mathcal{K}^{{\rm res}}$, which led to an explicit expression of the quotient $\kappa_{s_\varphi w}^\mathcal{O}/\kappa_w^\mathcal{O}$ (see~\cite[Lemma 5.6\,(1)]{SSV2}, as well as case 2 of the proof of~\cite[Lemma 5.9]{SSV2}).
This can be easily extended to $\kappa_w^\mathcal{O}$ for arbitrary parameters $\mathbf{k}\in\mathcal{K}$. It leads to the formula
\begin{gather*}
\frac{\kappa_{s_\varphi w}^\mathcal{O}}{\kappa_w^\mathcal{O}}=\mathbf{k}_\varphi^{-\chi(w^{-1}\varphi)\chi_e(\varphi(wc^\mathcal{O}))}
\mathbf{k}_{\frac{\varphi}{2}}^{\chi(w^{-1}\varphi)\chi_o(\varphi(wc^\mathcal{O}))}
\prod_{\alpha\in\Phi_0^+}\mathbf{k}_\alpha^{\chi_e(\alpha(c^\mathcal{O}))\alpha(w^{-1}\varphi^\vee)}\mathbf{k}_{\frac{\alpha}{2}}^{-\chi_o(\alpha(c^\mathcal{O}))\alpha(w^{-1}\varphi^\vee)},
\end{gather*}
where $\chi$ is given by~\eqref{chi}. Hence
\begin{gather*}
{\rm coeff}_w^\mathcal{O}=k_r^{-\chi(w^{-1}\varphi)\chi_e(\varphi(wc^\mathcal{O}))}u_r^{\chi(w^{-1}\varphi)\chi_o(\varphi(wc^\mathcal{O}))}
\biggl(\mathfrak{s}\prod_{\alpha\in\Phi_0^+}\mathbf{k}_\alpha^{\chi_e(\alpha(c^\mathcal{O}))\alpha} \mathbf{k}_{\frac{\alpha}{2}}^{-\chi_o(\alpha(c^\mathcal{O}))\alpha}\biggr)^{w^{-1}\varphi^\vee}\\
\hphantom{{\rm coeff}_w^\mathcal{O}=}{}
+\bigl(k_0-k_0^{-1}\bigr)\chi_+\bigl(w^{-1}\varphi\bigr)\chi_e\bigl(\varphi\bigl(wc^\mathcal{O}\bigr)\bigr)+
\bigl(u_0-u_0^{-1}\bigr)\chi_-\bigl(w^{-1}\varphi\bigr)\chi_o\bigl(\varphi\bigl(wc^\mathcal{O}\bigr)\bigr),
\end{gather*}
where the factor in big brackets in the first line is considered as element in $\mathbf{T}$ with value at~${\lambda\in Q^\vee}$ given by
\[
\mathfrak{s}^\lambda\prod_{\alpha\in\Phi_0^+}\mathbf{k}_\alpha^{\chi_e(\alpha(c^\mathcal{O}))\alpha(\lambda)}\,\mathbf{k}_{\frac{\alpha}{2}}^{-\chi_o(\alpha(c^\mathcal{O}))\alpha(\lambda)}.
\]
Now we pick $\mathfrak{s}\in\mathbf{T}$ to be
\begin{equation}\label{s}
\mathfrak{s}_\mathcal{O}:=\prod_{\alpha\in\Phi_0^+}\bigl(\mathbf{k}_{\alpha}\mathbf{k}_{(\alpha,1)}\bigr)^{-\frac{\chi_e(\alpha(c^\mathcal{O}))}{2}\alpha}
\bigl(\mathbf{k}_{\frac{\alpha}{2}}\mathbf{k}_{(\frac{\alpha}{2},\frac{1}{2})}\bigr)^{\frac{\chi_o(\alpha(c^\mathcal{O}))}{2}\alpha},
\end{equation}
whose value at $\lambda\in Q^\vee$ is
\[
\mathfrak{s}^\lambda_\mathcal{O}:=\prod_{\alpha\in\Phi_0^+}\bigl(\mathbf{k}_{\alpha}\mathbf{k}_{(\alpha,1)}\bigr)^{-\frac{\chi_e(\alpha(c^\mathcal{O}))\alpha(\lambda)}{2}}
\bigl(\mathbf{k}_{\frac{\alpha}{2}}\mathbf{k}_{(\frac{\alpha}{2},\frac{1}{2})}\bigr)^{\frac{\chi_o(\alpha(c^\mathcal{O}))\alpha(\lambda)}{2}}.
\]
When \smash{$\alpha\bigl(Q^\vee\bigr)=\mathbb{Z}$}, the factor in this product should be read as
\[
\mathbf{k}_\alpha^{-\chi_e(\alpha(c^\mathcal{O}))\alpha(\lambda)}\,\mathbf{k}_{\frac{\alpha}{2}}^{\chi_o(\alpha(c^\mathcal{O}))\alpha(\lambda)}=
\mathbf{k}_{\alpha}^{(\chi_o(\alpha(c^\mathcal{O}))-\chi_e(\alpha(c^\mathcal{O})))\alpha(\lambda)}
\]
\big(recall that \smash{$\mathbf{k}_\alpha=\mathbf{k}_{(\alpha,1)}=\mathbf{k}_{\frac{\alpha}{2}}=
\mathbf{k}_{(\frac{\alpha}{2},\frac{1}{2})}$} when $\alpha\smash{\bigl(Q^\vee\bigr)}=\mathbb{Z}$\big).

For the remainder of the proof, we set $\mathfrak{s}=\mathfrak{s}_\mathcal{O}$, hence the linear map $\psi_t^\mathcal{O}\colon \mathbb{H}\twoheadrightarrow\mathbf{F}[\mathcal{O}]$ is now given by
\begin{equation}\label{psitO}
\psi_t^\mathcal{O}\bigl(x^\lambda T_wY^\mu\bigr)=\kappa_w^\mathcal{O}(\mathfrak{s}_\mathcal{O}t)^{-\mu}x^{\lambda+wc^\mathcal{O}}
\end{equation}
for $\lambda,\mu\in Q^\vee$ and $w\in W_0$.
We get
\begin{gather}
{\rm coeff}_w^\mathcal{O}=k_r^{-\chi(w^{-1}\varphi)\chi_e(\varphi(wc^\mathcal{O}))}
u_r^{\chi(w^{-1}\varphi)\chi_o(\varphi(wc^\mathcal{O}))}\nonumber\\
\hphantom{{\rm coeff}_w^\mathcal{O}=}\quad{}
 \times\prod_{\alpha\in\Phi_0^+: \alpha\smash{(Q^\vee)}=\mathbb{Z}_e}
\bigl(\mathbf{k}_\alpha \mathbf{k}_{(\alpha,1)}^{-1}\bigr)^{\frac{\chi_e(\alpha(c^\mathcal{O}))\alpha(w^{-1}\varphi^\vee)}{2}} \bigl(\mathbf{k}_{\frac{\alpha}{2}}^{-1}\mathbf{k}_{(\frac{\alpha}{2},\frac{1}{2})}\bigr)^{\frac{\chi_o(\alpha(c^\mathcal{O}))\alpha(w^{-1}\varphi^\vee)}{2}}\nonumber\\
\hphantom{{\rm coeff}_w^\mathcal{O}=}{}
+\bigl(k_0-k_0^{-1}\bigr)\chi_+\bigl(w^{-1}\varphi\bigr)\chi_e\bigl(\varphi\bigl(wc^\mathcal{O}\bigr)\bigr)+
\bigl(u_0-u_0^{-1}\bigr)\chi_-\bigl(w^{-1}\varphi\bigr)\chi_o\bigl(\varphi\bigl(wc^\mathcal{O}\bigr)\bigr).
\label{coeffmid}
\end{gather}
Recall that there only exist roots $\alpha\in\Phi_0$ with $\alpha\smash{\bigl(Q^\vee\bigr)}=\mathbb{Z}_e$ when $\Phi_0$ is of type ${\rm C}_r$. In this case
\smash{$\bigl\{\alpha\in\Phi_0\mid \alpha\bigl(Q^\vee\bigr)=\mathbb{Z}_e\bigr\}$} is the set $\Phi_{0,\ell}$ of long roots in $\Phi_0$, and
\[
\bigl\{\alpha\in\Phi_{0,\ell} \mid \alpha\bigl(w^{-1}\varphi^\vee\bigr)\not=0\bigr\}=\big\{w^{-1}\varphi,-w^{-1}\varphi\big\}.
\]
We thus have
\begin{gather*}
\prod_{\alpha\in\Phi_0^+\colon \alpha\smash{(Q^\vee)}=\mathbb{Z}_e}
\bigl(\mathbf{k}_\alpha \mathbf{k}_{(\alpha,1)}^{-1}\bigr)^{\frac{\chi_e(\alpha(c^\mathcal{O}))\alpha(w^{-1}\varphi^\vee)}{2}} \bigl(\mathbf{k}_{\frac{\alpha}{2}}^{-1}\mathbf{k}_{(\frac{\alpha}{2},\frac{1}{2})}\bigr)^{\frac{\chi_o(\alpha(c^\mathcal{O}))\alpha(w^{-1}\varphi^\vee)}{2}}\\
\qquad{}=
\bigl(k_0^{-1}k_r\bigr)^{\chi(w^{-1}\varphi)\chi_e(\varphi(wc^\mathcal{O}))}
\bigl(u_0u_r^{-1}\bigr)^{\chi(w^{-1}\varphi)\chi_o(\varphi(wc^\mathcal{O}))}.\label{hhhelp}
\end{gather*}
Note that this formula is correct for $\Phi_0$ of arbitrary type. Indeed, if $\Phi_0$ is not of type $C_r$ then the product on the left-hand side is an empty product, and the right-hand side also reduces to~$1$.
Substituting~\eqref{hhhelp} into
\eqref{coeffmid}, we see that
the dependence on $k_r$ and $u_r$ drops out, and we end up with the formula
\begin{gather*}
{\rm coeff}_w^\mathcal{O}=k_0^{-\chi(w^{-1}\varphi)\chi_e(\varphi(wc^\mathcal{O}))}
u_0^{\chi(w^{-1}\varphi)\chi_o(\varphi(wc^\mathcal{O}))}\\
\hphantom{{\rm coeff}_w^\mathcal{O}=}{}
+\bigl(k_0-k_0^{-1}\bigr)\chi_+\bigl(w^{-1}\varphi\bigr)\chi_e\bigl(\varphi\bigl(wc^\mathcal{O}\bigr)\bigr)+
\bigl(u_0-u_0^{-1}\bigr)\chi_-\bigl(w^{-1}\varphi\bigr)\chi_o\bigl(\varphi\bigl(wc^\mathcal{O}\bigr)\bigr).
\end{gather*}
Now note that
\begin{gather*}
k_0^{-\chi(w^{-1}\varphi)\chi_e(\varphi(wc^\mathcal{O}))}
u_0^{\chi(w^{-1}\varphi)\chi_o(\varphi(wc^\mathcal{O}))}=\bigl(k_0\chi_-\bigl(w^{-1}\varphi\bigr)+k_0^{-1}\chi_+\bigl(w^{-1}\varphi\bigr)\bigr)\chi_e \bigl(\varphi\bigl(wc^\mathcal{O}\bigr)\bigr)\\
\hphantom{k_0^{-\chi(w^{-1}\varphi)\chi_e(\varphi(wc^\mathcal{O}))}
u_0^{\chi(w^{-1}\varphi)\chi_o(\varphi(wc^\mathcal{O}))}=}{}
+\bigl(u_0\chi_+\bigl(w^{-1}\varphi\bigr)+u_0^{-1}\chi_-\bigl(w^{-1}\varphi\bigr)\bigr)\chi_o\bigl(\varphi\bigl(wc^\mathcal{O}\bigr)\bigr),
\end{gather*}
and hence
\begin{equation}\label{coefffin}
{\rm coeff}_w^\mathcal{O}=k_0\chi_e\bigl(\varphi\bigl(wc^\mathcal{O}\bigr)\bigr)+u_0\chi_o\bigl(\varphi\bigl(wc^\mathcal{O}\bigr)\bigr)=
k_0^{\chi_e(\varphi(wc^\mathcal{O}))}u_0^{\chi_o(\varphi(wc^\mathcal{O}))}.
\end{equation}
Note that \smash{$y\mapsto k_0^{\chi_e(\varphi(y))}u_0^{\chi_o(\varphi(y))}$} is $\tau\smash{\bigl(Q^\vee\bigr)}$-invariant (compare with the proof of Theorem~\ref{theoremH}),
hence~\eqref{formula0} and~\eqref{coefffin} lead to the formula
\begin{equation*}
\frac{(\mathfrak{s}_\mathcal{O}t)^\mu}{\kappa_w^\mathcal{O}}\psi_{t}^\mathcal{O}\bigl(T_0x^\lambda T_wY^\mu\bigr)=\mathcal{D}_0(x^y)
\end{equation*}
for $y=\lambda+wc^\mathcal{O}\in\mathcal{O}$ and $\mu\in Q^\vee$, where $\mathcal{D}_0$ is the linear operator on $\mathbf{F}[\mathcal{O}]$ defined by
\[
\mathcal{D}_0(x^y):=
k_0^{\chi_e(\overline{\alpha}_0(y))}u_0^{\chi_o(\overline{\alpha}_0(y))}s_{0,t}x^y+\bigl(k_0-k_0^{-1}\bigr)\nabla_0^e(x^y)+\bigl(u_0-u_0^{-1}\bigr)\nabla_0^o(x^y)
\]
for $y\in\mathcal{O}$.

In conclusion, the kernel of $\psi_t^\mathcal{O}\colon \mathbb{H}\twoheadrightarrow\mathbf{F}[\mathcal{O}]$ (see~\eqref{psitO})
is a left $\mathbb{H}$-module, and the resulting isomorphism \smash{$\mathbb{H}/\textup{ker}\bigl(\psi_t^\mathcal{O}\bigr)\overset{\sim}{\longrightarrow}\mathbf{F}[\mathcal{O}]$}
extends the quasi-polynomial $H^X$-action $\pi$ on $\mathbf{F}[\mathcal{O}]$ to an action of $\mathbb{H}$ with $T_0\in H\subset\mathbb{H}$ acting by
$\mathcal{D}_0$, hence the corresponding representation map is \smash{$\pi_t^\mathcal{O}$}. This concludes the proof.
\end{proof}
